%% =====================================================================
%%  Capítulo 10 -- Respuesta frecuencial y diagramas de Bode
%%  PILOTO del formato nuevo (noviembre de 2026).
%%  Integra: Handout8 (versión anterior), notas de la Clase 1, las cuatro
%%  presentaciones Bode1-Bode4 (v3) y los códigos de Bode_4Clases/matlab.
%%  Correcciones respecto al handout anterior (ver \S\ref{sec:rf-cierre}):
%%    adelanto w_M = 8.85 (no 14); atraso con holgura w = 1.68, PM = 44.8°;
%%    rango de Kc: GM = +40 dB (no -40 dB); w_GM asintótico 31.6 (no 50);
%%    Kc exacto 33 056; se usa M_r (no M_P) para el pico de resonancia.
%%  Sin retardos (decisión del curso): el PM se explica como presupuesto de
%%  fase frente a dinámicas no modeladas (un sensor de primer orden).
%% =====================================================================

\capitulo[cap:rf]{Respuesta frecuencial y diagramas de Bode}
  {Compramos libros porque creemos que compramos el tiempo para leerlos.}
  {Warren Zevon}

\begin{objetivos}[partes/10-respuesta-frecuencial.html]
  \item Explicar por qué un sistema lineal y estable responde a una senoide con otra senoide, y calcular su amplitud y su fase.
  \item Dibujar a mano el diagrama de Bode asintótico de una malla abierta con ganancia, integradores y polos y ceros reales, y estimar cuánto se equivoca.
  \item Leer en un diagrama de Bode los márgenes de ganancia y de fase, decidir si el lazo cerrado es estable y cuánta ganancia admite.
  \item Predecir el sobrepico, la velocidad y el error en estado estacionario del lazo cerrado mirando solamente la malla abierta.
  \item Diseñar un controlador proporcional, un compensador en adelanto, en atraso o atraso-adelanto para cumplir especificaciones de error y de margen de fase, y verificar el diseño en \matlab/\simulink.
\end{objetivos}

%% ---------------------------------------------------------------------
\section{¿Por qué la frecuencia?}
\label{sec:rf-porque}
%% ---------------------------------------------------------------------

Hasta ahora hemos analizado los sistemas de control de dos maneras: con la función de transferencia, mirando dónde quedan los polos en el plano-$s$, y con la representación de estado, mirando los valores propios de $\Am$. Las dos exigen \emph{tener el modelo}. La respuesta frecuencial cambia la pregunta. En lugar de preguntar dónde están los polos, pregunta
\begin{center}
\sffamily\large\color{ctrlazul}¿qué le hace el sistema a una senoide?
\end{center}
La respuesta, como veremos en la \cref{sec:rf-senoide}, son dos números por cada frecuencia: cuánto cambia la amplitud y cuánto se corre la fase. Lo valioso es que esos dos números se pueden \emph{medir} en un laboratorio con un generador de señales y un osciloscopio, aun cuando nadie haya escrito nunca la función de transferencia. Y, sobre todo, con esos dos números dibujados en una sola gráfica, el \clave{diagrama de Bode}, vamos a poder responder las tres preguntas del curso (estabilidad, respuesta transitoria y error) y además \emph{diseñar}: la \cref{fig:rf-mapa} es el mapa del capítulo.

\begin{figure}[ht]
\centering
\begin{tikzpicture}[font=\sffamily\small,
  caja/.style={draw=#1,fill=#1!8,thick,rounded corners=2pt,align=center,minimum height=10mm,text width=27mm},
  flecha/.style={-{Stealth[length=2.2mm]},thick,draw=ctrlgris}]
  \node[caja=ctrlnoche] (sen) {una senoide entra\\ y sale otra};
  \node[caja=ctrlazul,right=6mm of sen] (bode) {diagrama de Bode\\ de $\Gma(\jw)$};
  \node[caja=ctrlverde,right=9mm of bode,yshift=13mm] (est) {estabilidad\\ \footnotesize $\GM$, $\PM$ (\S\ref{sec:rf-estabilidad})};
  \node[caja=ctrlvioleta,right=9mm of bode] (tra) {transitorio\\ \footnotesize $\PM\to\zeta$, $\wbw$ (\S\ref{sec:rf-lazocerrado})};
  \node[caja=ctrlambar,right=9mm of bode,yshift=-13mm] (err) {error\\ \footnotesize tipo, $\kp,\kv,\ka$ (\S\ref{sec:rf-tipo})};
  \node[caja=ctrlcoral,below=7mm of bode] (dis) {diseño: $\Kc$, adelanto,\\ atraso (\S\ref{sec:rf-ajuste}--\ref{sec:rf-pid})};
  \draw[flecha] (sen) -- (bode);
  \draw[flecha] (bode.east) -- (est.west);
  \draw[flecha] (bode.east) -- (tra.west);
  \draw[flecha] (bode.east) -- (err.west);
  \draw[flecha] (bode) -- (dis);
\end{tikzpicture}
\caption{Mapa del capítulo. Todo sale de un mismo dibujo: la malla abierta $\Gma(\jw)$ en magnitud y fase.}
\label{fig:rf-mapa}
\end{figure}

%% ---------------------------------------------------------------------
\subsection{Un poco de historia: el control nace en un laboratorio de teléfonos}

En los años veinte, las líneas telefónicas transcontinentales de AT\&T necesitaban cientos de amplificadores de tubos de vacío en cascada. La distorsión de cada uno se acumulaba y la ganancia derivaba con la temperatura y con el envejecimiento de los tubos. Harold S.\ Black propuso en 1927 una solución que hoy parece obvia: \emph{retroalimentar} negativamente el amplificador, sacrificando ganancia a cambio de precisión y linealidad~\citep{black34}.\nota{Black cuenta que la idea se le ocurrió el 2 de agosto de 1927 en el ferry de Lackawanna, camino al trabajo, y que la anotó en una página del \emph{New York Times}~\citep{black77}. La patente, presentada en 1928, se concedió en 1937~\citep{blackpat}.} El precio fue un problema nuevo: con mucha ganancia de lazo, los amplificadores \emph{oscilaban}.

\notamargen{\textbf{Pregunta para la clase.} ¿Por qué una compañía telefónica necesitaba una teoría de estabilidad?}
Harry Nyquist dio en 1932 el criterio para predecir la estabilidad del lazo cerrado a partir de mediciones de la malla abierta~\citep{Nyq32}, y Hendrik W.\ Bode, entre 1940 y 1945, convirtió esas ideas en un método de diseño de regla y lápiz: las gráficas de magnitud en decibeles y de fase contra el logaritmo de la frecuencia que hoy llevan su nombre, las aproximaciones asintóticas, los márgenes de ganancia y de fase, y la relación entre la pendiente de la magnitud y la fase en los sistemas de fase mínima~\citep{Bod40,bode45}. Durante la Segunda Guerra Mundial estos métodos salieron de los teléfonos y llegaron a los servomecanismos de radares y cañones~\citep{JaNiPh47}; el lugar de las raíces de Evans llegaría después~\citep{Eva50}. La \cref{fig:rf-linea} resume la cronología.

\begin{figure}[ht]
\centering
\includegraphics[width=\linewidth]{cap10/f_linea_tiempo}
\caption{De Black a Evans: casi todo el control clásico en frecuencia salió de Bell Labs (en ámbar, los aportes de Bode).}
\label{fig:rf-linea}
\end{figure}

\begin{historia}[Hendrik Wade Bode (1905--1982)]
Ingeniero de redes y amplificadores en Bell Labs. Su problema era diseñar retroalimentación con mucha ganancia y sin oscilaciones. En 1940 mostró que en los sistemas de fase mínima la magnitud determina la fase: basta mirar la pendiente. En la guerra trabajó en directores de tiro antiaéreo y en 1967 pasó a Harvard como profesor de ingeniería de sistemas. Su gran aporte práctico fueron las aproximaciones con líneas rectas, que permiten dibujar a mano, en minutos, el comportamiento de un sistema en \emph{todas} las frecuencias.
\end{historia}

%% ---------------------------------------------------------------------
\subsection{El lenguaje: malla abierta, lazo abierto y \emph{open loop}}
\label{sec:rf-lenguaje}

En el curso distinguimos dos funciones que en los libros en inglés suelen aparecer con el mismo nombre. Considere el lazo de la \cref{fig:rf-mallalazo}(a), con camino directo $G(s)$ y retroalimentación $\Hs(s)$.

\begin{figure}[ht]
\centering
\includegraphics[width=\linewidth]{cap10/malla_vs_lazo}
\caption{(a) La malla abierta es lo que recorre una señal al dar una vuelta completa al lazo. (b) El lazo abierto aparece al reescribir el mismo lazo con retroalimentación unitaria.}
\label{fig:rf-mallalazo}
\end{figure}

\begin{definicion}[label=def:rf-malla]{Malla abierta y lazo abierto}
La \clave{malla abierta} es el producto de todo lo que recorre una señal al dar una vuelta completa al lazo, \emph{sin} el signo del comparador:
\begin{equation}
  \Gma(s)=G(s)\,\Hs(s).
  \label{eq:rf-mallaabierta}
\end{equation}
Es la función con la que se estudia la estabilidad del lazo cerrado y la que se dibuja en el diagrama de Bode. Si el lazo tiene válvula, planta, sensor y transmisor, $\Gma=\Gv\,\Gp\,\Gs\,\Gt$.

El \clave{lazo abierto} es la función que relaciona el error con la salida cuando el lazo se reescribe con retroalimentación unitaria:
\begin{equation}
  \Gla(s)=\frac{G(s)}{1+G(s)\Hs(s)-G(s)}.
  \label{eq:rf-lazoabierto}
\end{equation}
Es la que se usa para el tipo del sistema y los errores en estado estacionario (capítulo de error).
\end{definicion}

Las dos llevan a la misma ecuación característica,
\begin{equation}
  1+G(s)\Hs(s)=0\quad\Longleftrightarrow\quad 1+\Gla(s)=0,
  \label{eq:rf-ecuacioncaracteristica}
\end{equation}
y cuando $\Hs(s)=1$ (el caso de casi todos nuestros ejemplos) coinciden: $\Gma(s)=\Gla(s)=G(s)$. Por eso los textos las mezclan. Ogata, por ejemplo, llama \emph{open-loop transfer function} a $B(s)/E(s)=G(s)H(s)$, que es nuestra malla abierta, y \emph{feedforward transfer function} a $C(s)/E(s)=G(s)$~\citep{ogata10}; Åström y Murray la llaman \emph{loop transfer function}, $L=PC$~\citep{AsMu10}; en electrónica se dice \emph{loop gain}, $A\beta$. La regla para no confundirse es preguntar \emph{qué se corta y qué se mide}.

\begin{cuidado}
Con $\Gma(s)$ estudiamos estabilidad y diseñamos en frecuencia; con $\Gla(s)$ calculamos el tipo y el error. Cuando $\Hs\neq1$ son funciones \emph{distintas}: un sensor con dinámica cambia la malla abierta y, con ella, los márgenes, aunque el error se calcule con otra función.
\end{cuidado}

\paragraph{Cortar el lazo}
Imagine que corta el lazo justo después del comparador (\cref{fig:rf-corte}), inyecta una señal $x$ en el corte y mide lo que regresa, $z$. La señal recorre $G$, luego $\Hs$, y llega al comparador, que le cambia el signo:
\begin{equation}
  \frac{Z(s)}{X(s)}=-\,G(s)\,\Hs(s)=-\,\Gma(s).
\end{equation}
Si para alguna frecuencia $\omega$ lo que vuelve es exactamente $-x$ (misma amplitud, invertido), es decir, si $\Gma(\jw)=-1$, el comparador lo convierte en $+x$: el lazo se alimenta solo y oscila sin necesidad de entrada. Por eso el punto crítico del análisis en frecuencia es $-1$ (magnitud $1=0\dB$ y fase $-180\grados$), y por eso el signo del comparador \emph{no} se incluye en $\Gma(s)$. Bode llamó \clave{diferencia de retorno} a $1+\Gma(s)$: la diferencia entre la señal inyectada y la que regresa. Sus ceros son los polos del lazo cerrado,
\begin{equation}
  \Tlc(s)=\frac{Y(s)}{R(s)}=\frac{G(s)}{1+\Gma(s)}.
\end{equation}

\begin{figure}[ht]
\centering
\includegraphics[width=0.62\linewidth]{cap10/cortar_lazo}
\caption{Cortar el lazo: se inyecta $x$ y se mide lo que regresa, $z=-\Gma x$.}
\label{fig:rf-corte}
\end{figure}

\paragraph{El mismo lenguaje en electrónica}
\notamargen{En las hojas de datos de los amplificadores operacionales, la \emph{open-loop gain} es $A$, del orden de $10^5$.}
En el amplificador de Black (\cref{fig:rf-black}), $A$ es la ganancia del amplificador sin retroalimentar y $\beta$ la fracción de la salida que se resta a la entrada:
\begin{equation}
  A_f=\frac{A}{1+A\beta}\ \xrightarrow{\ A\beta\gg 1\ }\ \frac{1}{\beta}.
\end{equation}
Con $A=10^5$ y $\beta=0.01$ resulta $A_f=99.90$; si los tubos envejecen y $A$ cae a la mitad, $A_f=99.80$. La ganancia del amplificador cayó un $50\%$ y la del amplificador retroalimentado apenas un $0.1\%$. Es la misma razón por la que retroalimentamos plantas: la ganancia de lazo $A\beta\leftrightarrow\Gma(s)$ compra insensibilidad a las variaciones del sistema.

\begin{figure}[ht]
\centering
\includegraphics[width=0.5\linewidth]{cap10/black}
\caption{El amplificador retroalimentado de Black: $A\beta$ es la malla abierta.}
\label{fig:rf-black}
\end{figure}

%% ---------------------------------------------------------------------
\section{Una senoide entra, una senoide sale}
\label{sec:rf-senoide}
%% ---------------------------------------------------------------------

\begin{definicion}[label=def:rf-respuesta]{Respuesta en frecuencia}
La \clave{respuesta en frecuencia} de un sistema lineal, invariante en el tiempo y \emph{estable} es su respuesta en estado estacionario a una entrada senoidal. La salida es una senoide de la \emph{misma} frecuencia; solo cambian la amplitud y la fase.
\end{definicion}

El supuesto de estabilidad (en malla abierta, cuando lo que se excita es $\Gma(s)$) es esencial: si hay polos en el semiplano derecho, los transitorios no se apagan y no hay un estado estacionario que medir. Lo mantendremos durante todo el capítulo.

\subsection{¿Por qué sale una senoide? Deducción}
\label{sec:rf-deduccion}

Sea $x(t)=x_0\sin(\omega t)$, cuya transformada es
\begin{equation}
  X(s)=\frac{x_0\,\omega}{s^2+\omega^2}=\frac{x_0\,\omega}{(s-\jw)(s+\jw)}.
\end{equation}
Si $\Gma(s)$ tiene polos $-p_1,\dots,-p_n$ (distintos, para simplificar), la salida se expande en fracciones parciales como
\begin{equation}
  Z(s)=\Gma(s)X(s)=\underbrace{\frac{K_1}{s+p_1}+\cdots+\frac{K_n}{s+p_n}}_{\text{modos del sistema}}
  +\underbrace{\frac{a}{s-\jw}+\frac{\bar a}{s+\jw}}_{\text{modo de la entrada}}.
\end{equation}
En el tiempo, los términos $K_ie^{-p_it}$ se extinguen si todos los polos tienen parte real negativa. Lo que queda es el modo de la entrada. El residuo en $s=\jw$ es
\begin{equation}
  a=\Big[(s-\jw)\,\Gma(s)\,\frac{x_0\,\omega}{(s-\jw)(s+\jw)}\Big]_{s=\jw}
  =\Gma(\jw)\,\frac{x_0}{2j},
\end{equation}
y el de $s=-\jw$ es su conjugado. Escribiendo $\Gma(\jw)=|\Gma(\jw)|\,e^{j\phi}$, con $\phi=\angle\Gma(\jw)$, y usando que $\Gma(-\jw)$ es el conjugado de $\Gma(\jw)$ (coeficientes reales):
\begin{equation}
  z_{ss}(t)=a\,e^{\jw t}+\bar a\,e^{-\jw t}
  =x_0\,|\Gma(\jw)|\,\frac{e^{j(\omega t+\phi)}-e^{-j(\omega t+\phi)}}{2j}
  =x_0\,|\Gma(\jw)|\,\sin(\omega t+\phi).
  \label{eq:rf-zss}
\end{equation}

\begin{ideaclave}
Para cada frecuencia, el sistema \emph{multiplica} la amplitud por $|\Gma(\jw)|$ y \emph{desplaza} la fase en $\phi=\angle\Gma(\jw)$. El diagrama de Bode es la gráfica de esos dos números en función de $\omega$.
\end{ideaclave}

\begin{ejemplo}[label=ej:rf-seno]{Un primer orden excitado con $\sin t$}
Calcule la salida en estado estacionario de $\Gma(s)=1/(s+10)$ cuando $x(t)=\sin t$.
\solucion Con $\omega=1$\,rad/s,
\[
  \Gma(j1)=\frac{1}{10+j}\ \Rightarrow\ |\Gma(j1)|=\frac{1}{\sqrt{101}}=0.0995,\qquad
  \phi=-\tan^{-1}\frac{1}{10}=-5.71\grados,
\]
y por \eqref{eq:rf-zss}, $z_{ss}(t)=0.0995\,\sin(t-5.71\grados)$. La salida es diez veces más pequeña y llega $5.7\grados$ tarde. Eso es \emph{todo} lo que el diagrama de Bode guarda para $\omega=1$. La \cref{fig:rf-seno} muestra que la simulación completa coincide con esta predicción en cuanto se apaga el transitorio $e^{-10t}$ (unas cuatro constantes de tiempo, $0.4$\,s).
\end{ejemplo}

\begin{figure}[ht]
\centering
\includegraphics[width=\linewidth]{cap10/f_seno_primer_orden}
\caption{$1/(s+10)$ excitado con $\sin t$: la simulación completa (azul) y la predicción de estado estacionario \eqref{eq:rf-zss} (coral) coinciden después del transitorio.}
\label{fig:rf-seno}
\end{figure}

\script[lsim y la fórmula dan lo mismo]{cap10_respuesta_frecuencial/Clase1_RespuestaFrecuencial/C1_Ej1_RespuestaSenoidal.m}
\codigorepo{cap10_respuesta_frecuencial/Clase1_RespuestaFrecuencial/C1_Ej1_RespuestaSenoidal.m}{senoide}{ejemplo \ref{ej:rf-seno}}

\video{V1_maquina_bode}{La máquina de Bode: se barre la frecuencia, se miden los dos números y el Bode se dibuja punto a punto.}
\begin{modelo}[RespuestaSenoidal.slx]
El script \cmd{crearModelo\_RespuestaSenoidal.m} construye por código el diagrama \emph{Sine Wave} $\to$ $G(s)$ $\to$ \emph{Scope}, lo simula y mide automáticamente la amplitud y el desfase en el último periodo, comparándolos con \cmd{bode(G,w)}. Cambie \cmd{w}, \cmd{num} y \cmd{den} en el \emph{workspace} y repita: está midiendo el diagrama de Bode como se hacía en Bell Labs, un punto a la vez.
\end{modelo}

\begin{pausa}
Sin calcular nada: si la entrada de $1/(s+10)$ fuera $\sin(100t)$, ¿la salida sería más grande o más pequeña que la del \cref{ej:rf-seno}? ¿Llegaría más o menos atrasada?
\respuesta{Más pequeña y más atrasada: $|\Gma(j100)|=1/\sqrt{10\,100}=0.00995$ (diez veces menor que en $\omega=1$) y $\phi=-\tan^{-1}10=-84.3\grados$. Un polo real atenúa y atrasa más a medida que la frecuencia sube.}
\end{pausa}

\subsection{¿Qué ganamos trabajando en frecuencia?}
\begin{enumerate}
  \item \textbf{Se mide.} Hay generadores de senoides en rangos enormes de frecuencia y amplitud: el modelo se puede obtener experimentalmente.
  \item \textbf{Se multiplica sumando.} En escala logarítmica, los bloques en cascada se suman (\cref{sec:rf-decibel}).
  \item \textbf{Se diseña a ojo.} La estabilidad relativa (márgenes), el error en estado estacionario y la velocidad de respuesta se leen en el mismo diagrama.
  \item \textbf{Se filtra.} El ruido de medición y las perturbaciones suelen vivir en bandas de frecuencia conocidas; el diseñador ve de inmediato qué deja pasar el lazo.
\end{enumerate}

%% ---------------------------------------------------------------------
\section{El decibel}
\label{sec:rf-decibel}
%% ---------------------------------------------------------------------

\subsection{Un invento de ingenieros telefónicos}

Antes de 1923, la pérdida de una línea telefónica se expresaba en \emph{millas de cable estándar} (MSC): cuántas millas de un cable de referencia atenuaban lo mismo a $5000$\,rad/s ($\approx796$\,Hz, una frecuencia representativa de la voz). La unidad dependía de la frecuencia y del cable. En 1924, W.~H.~Martin, de Bell Labs, propuso reemplazarla por la \emph{transmission unit} (TU): diez veces el logaritmo decimal de la razón de potencias~\citep{martin24}. La TU casi coincidía con la milla de cable,
\begin{equation}
  1\ \text{MSC}\approx\frac{1}{1.056}\ \text{TU}\approx 0.95\ \text{dB},
\end{equation}
de modo que los ingenieros no tuvieron que cambiar su intuición. En 1928 el Bell System adoptó el nombre \emph{decibel}, en honor a Alexander Graham Bell, siguiendo la propuesta internacional de llamar \emph{bel} a la unidad decimal y \emph{neper} (por John Napier) a la natural~\citep{martin29}:
\begin{gather}
  \text{bel}=\log_{10}\frac{P_2}{P_1},\qquad
  \text{dB}=10\log_{10}\frac{P_2}{P_1},\qquad
  1\ \text{Np}=20\log_{10}e=8.686\ \text{dB}.
\end{gather}
El bel resultó demasiado grande para el trabajo diario; sobrevivió su décima parte.

\subsection{¿Por qué 10 log y por qué 20 log?}

El decibel se define para \emph{potencias}. Si la potencia se disipa sobre la misma resistencia, $P=V^2/R$ y
\begin{equation}
  10\log_{10}\frac{V_2^2/R}{V_1^2/R}=20\log_{10}\frac{V_2}{V_1}.
\end{equation}
En el diagrama de Bode $|\Gma(\jw)|$ es una razón de \emph{amplitudes} (salida sobre entrada); por eso siempre se usa
\begin{equation}
  |\Gma|_{\mathrm{dB}}=20\log_{10}|\Gma(\jw)|.
\end{equation}

\begin{table}[ht]
\centering\small
\begin{tabular}{@{}cc l@{}}
\toprule
$|G|$ & dB & para recordar\\
\midrule
$10$ & $+20$ & una década de amplitud\\
$2$ & $+6.02$ & el doble\\
$\sqrt2$ & $+3.01$ & la esquina de un cero real\\
$1$ & $0$ & cruce de ganancia\\
$1/\sqrt2$ & $-3.01$ & ancho de banda; esquina de un polo real\\
$0.1$ & $-20$ & una década menos\\
\bottomrule
\end{tabular}
\caption{Valores que conviene memorizar. La mitad de la \emph{potencia} es $-3$\,dB: por eso la frecuencia de corte de un filtro se define donde la magnitud cae $3$\,dB.}
\label{tab:rf-db}
\end{table}

\paragraph{Tres razones para usar logaritmos}
\begin{enumerate}
  \item \textbf{Cascadas.} Las ganancias se multiplican; en dB se suman, y las fases también se suman:
  \[
    20\log|G_1G_2G_3|=20\log|G_1|+20\log|G_2|+20\log|G_3|,\qquad
    \angle(G_1G_2G_3)=\angle G_1+\angle G_2+\angle G_3 .
  \]
  Por ejemplo, $\times10\cdot\times0.5\cdot\times4=\times20$ en amplitud equivale a $20-6.02+12.04=26.02$\,dB.
  \item \textbf{Rangos enormes.} Una razón de amplitudes de $10^6$ son ``solo'' $120$\,dB.
  \item \textbf{Percepción.} La sonoridad que percibimos crece aproximadamente con el logaritmo de la intensidad; un cambio de $1$\,dB apenas se nota.
\end{enumerate}

\notamargen{Cuando se agrega una referencia, el decibel deja de ser una razón y pasa a ser un \emph{nivel}: dBm ($1$\,mW), dBV ($1$\,V), dB SPL ($20\,\mu$Pa). Un número en dB sin referencia es siempre una razón.}

\begin{pausa}
Un amplificador de $+26$\,dB se conecta en cascada con un filtro que, a cierta frecuencia, atenúa $-32$\,dB. ¿Por cuánto queda multiplicada la amplitud de una senoide de esa frecuencia?
\respuesta{$26-32=-6$\,dB, es decir, $10^{-6/20}=0.50$: la amplitud queda a la mitad. En dB las cascadas se suman.}
\end{pausa}

%% ---------------------------------------------------------------------
\section{La forma de Bode y los cuatro factores}
\label{sec:rf-factores}
%% ---------------------------------------------------------------------

Toda función de transferencia racional se puede escribir en la \clave{forma de Bode}:
\begin{equation}
  \Gma(s)=\frac{\Kb\prod_{i=1}^{m}(1+\tau_i s)}
  {s^N\prod_{l=1}^{M}(1+\tau_l s)\prod_{k=1}^{n-M-N}\Big(\frac{s^2}{\omega_{n_k}^2}+\frac{2\zeta_k}{\omega_{n_k}}s+1\Big)},
  \label{eq:rf-formabode}
\end{equation}
con $m$ ceros reales, $N$ polos en el origen, $M$ polos reales y el resto en pares complejos conjugados; el grado relativo es $n-m$. (Si hay ceros complejos se tratan igual que los polos complejos, con el signo cambiado.) Lo importante es que cada factor está escrito como $(1+\tau s)$ y \emph{no} como $(s+a)$: así cada factor vale $1$ ($0\dB$) en bajas frecuencias y la ganancia $\Kb$ aparece sola. Tomando $20\log$ y el ángulo,
\begin{align}
  20\log|\Gma|&=20\log \Kb+\sum_i 20\log|1+\jw\tau_i|-20N\log\omega
  -\sum_l 20\log|1+\jw\tau_l|-\sum_k 20\log\Big|1-\tfrac{\omega^2}{\omega_{n_k}^2}+j\tfrac{2\zeta_k\omega}{\omega_{n_k}}\Big|,\notag\\
  \phi(\omega)&=\sum_i\tan^{-1}(\omega\tau_i)-N\cdot 90\grados
  -\sum_l\tan^{-1}(\omega\tau_l)
  -\sum_k\tan^{-1}\frac{2\zeta_k\omega_{n_k}\omega}{\omega_{n_k}^2-\omega^2}.
  \label{eq:rf-sumas}
\end{align}
El diagrama completo es la \emph{suma} de los diagramas de cada factor. Basta entonces aprender cuatro recetas.

%% ---------------------------------------------------------------------
\subsection{Factor 1: ganancia constante}
\[
  \Gma(s)=\Kb:\qquad 20\log|\Kb|\ \text{dB},\qquad
  \phi=\begin{cases}0\grados & \Kb>0,\\ \pm180\grados & \Kb<0.\end{cases}
\]
Una ganancia sube o baja \emph{toda} la curva de magnitud y no toca la fase. Esta idea, que parece trivial, es la base del diseño por ajuste de ganancia (\cref{sec:rf-ajuste}).

\subsection{Factor 2: polos y ceros en el origen}
\[
  \Gma(\jw)=\frac{1}{\jw}\ \Rightarrow\ 20\log\Big|\frac{1}{\jw}\Big|=-20\log\omega,\qquad\phi=-90\grados.
\]
En escala logarítmica es una recta de $-20\dBdec$ que pasa por $0\dB$ en $\omega=1$. En general $1/s^N$ da $-20N\dBdec$ y $-90N\grados$, y $s^N$ da $+20N\dBdec$ y $+90N\grados$. Todas estas rectas pasan por $0\dB$ en $\omega=1$ (\cref{fig:rf-simples}). El número de integradores, $N$, es el \emph{tipo} del sistema, y se lee en la pendiente de baja frecuencia (\cref{sec:rf-tipo}).

\begin{figure}[ht]
\centering
\includegraphics[width=\linewidth]{cap10/f_factores_simples}
\caption{Factores 1 y 2: una ganancia es una recta horizontal que no toca la fase; un integrador es una recta de $-20\dBdec$ que pasa por $0\dB$ en $\omega=1$ con fase constante de $-90\grados$.}
\label{fig:rf-simples}
\end{figure}

\subsection{Factor 3: polos y ceros reales}
Para el cero $\Gma(s)=s/\omega_1+1$,
\[
  20\log|\Gma(\jw)|=20\log\sqrt{\Big(\frac{\omega}{\omega_1}\Big)^2+1}\ \approx\
  \begin{cases}0\dB, & \omega\ll\omega_1,\\[2pt]
  20\log\dfrac{\omega}{\omega_1}\ \ (+20\dBdec), & \omega\gg\omega_1,\end{cases}
\]
y en la esquina $\omega=\omega_1$ el valor real es $20\log\sqrt2=3\dB$: el error máximo de la asíntota. La fase es $\phi=\tan^{-1}(\omega/\omega_1)$; su aproximación asintótica vale $0\grados$ hasta $\omega_1/10$, sube $45\grados$ por década pasando por $45\grados$ en $\omega_1$ y se queda en $90\grados$ desde $10\,\omega_1$ (el error en los extremos es $\tan^{-1}0.1=5.7\grados$). El polo es el espejo: $-20\dBdec$ desde $\omega_1$, $-3\dB$ en la esquina y fase de $0\grados$ a $-90\grados$ (\cref{fig:rf-polo}).

\begin{figure}[ht]
\centering
\includegraphics[width=0.88\linewidth]{cap10/f_polo_real}
\caption{Polo real: asíntotas (negro) y curva real (verde). El error máximo en magnitud es de $3\dB$ en la esquina; en fase, de $5.7\grados$ en $\omega_1/10$ y en $10\,\omega_1$.}
\label{fig:rf-polo}
\end{figure}

\begin{cuidado}[La trampa clásica: la ganancia escondida]
La regla solo vale si el factor está escrito como $s/\omega_1+1$. Como $s+10=10\,(s/10+1)$, el factor $1/(s+10)$ trae una ganancia $1/10=-20\dB$ escondida. Si se olvida sacarla, toda la curva queda corrida $20\dB$. Es el error más frecuente en las evaluaciones.
\end{cuidado}

\subsection{Factor 4: polos complejos (lo esencial)}
\label{sec:rf-complejos}
Para el factor normalizado de segundo orden,
\begin{equation}
  \Gma(s)=\frac{\wn^2}{s^2+2\zeta\wn s+\wn^2}\ \Rightarrow\
  |\Gma(\jw)|=\frac{\wn^2}{\sqrt{(\wn^2-\omega^2)^2+4\zeta^2\wn^2\omega^2}}.
\end{equation}
Las asíntotas son sencillas: $0\dB$ hasta $\wn$ y $-40\dBdec$ después; la fase va de $0\grados$ a $-180\grados$ y pasa por $-90\grados$ en $\wn$. Lo que la asíntota \emph{no} ve es el pico de resonancia. En $\omega=\wn$ el término $\wn^2-\omega^2$ se anula y
\begin{equation}
  \Gma(\jw_n)=\frac{1}{j2\zeta}\ \Rightarrow\ |\Gma(\jw_n)|=\frac{1}{2\zeta},\qquad
  \angle \Gma(\jw_n)=-90\grados,
  \label{eq:rf-pico}
\end{equation}
exactamente y para cualquier $\zeta$. Con $\zeta=0.1$ la curva real está $14\dB$ por encima de la asíntota (\cref{fig:rf-complejos}).

\begin{figure}[ht]
\centering
\includegraphics[width=0.88\linewidth]{cap10/f_complejos}
\caption{Polos complejos: magnitud y fase reales para varios $\zeta$, con las asíntotas en negro. Todas las fases pasan por $-90\grados$ en $\wn$.}
\label{fig:rf-complejos}
\end{figure}

\begin{ideaclave}
$\zeta$ pequeño $\Rightarrow$ pico grande en $\wn$ $\Rightarrow$ sobrepico grande en el tiempo. En el curso \alerta{no dibujamos asintóticos de segundo orden}: se dibujan las asíntotas, se marca el punto $20\log\frac{1}{2\zeta}$ en $\wn$ para recordar que la curva real se separa, y para la curva real se usa \cmd{bode}.
\end{ideaclave}

\video{V2_resonancia_zeta}{Tres vistas del mismo $\zeta$: los polos se acercan al eje imaginario, crece el pico del Bode y crece el sobrepico.}

\subsection{¿Por qué no basta la asíntota? Un ejemplo}
\label{sec:rf-sorpresa}

\begin{ejemplo}[label=ej:rf-sorpresa]{La asíntota dice ``estable''; la realidad, no}
¿Es estable el lazo cerrado con retroalimentación unitaria si
\begin{equation*}
  \Gma(s)=\frac{1000}{(s+1)(s^2+4s+100)}=\frac{10}{(s+1)\big(\frac{s^2}{100}+0.04s+1\big)}\ ?
\end{equation*}
\solucion Aquí $\wn=10$ y $\zeta=0.2$. Las asíntotas dicen: $20\dB$ hasta $\omega=1$, $-20\dBdec$ hasta $\omega=10$ (donde llegan exactamente a $0\dB$) y $-60\dBdec$ después. La fase real en $\omega=10$ es $-84.3\grados-90\grados=-174.3\grados$. Solo con asíntotas, el lazo cerrado parece estable por poco.

La curva real cuenta otra historia (\cref{fig:rf-sorpresa}): la resonancia suma $20\log\frac{1}{2(0.2)}=+7.96\dB$ en $\wn=10$, de modo que $|\Gma(j10)|=7.92\dB$. El cruce por $0\dB$ se corre a $12.7$\,rad/s, donde la fase ya pasó $-180\grados$: el margen de fase es $-45.5\grados$ y el lazo cerrado es \alerta{inestable}. Se comprueba con Routh-Hurwitz sobre $s^3+5s^2+104s+1100$: como $5\cdot104=520<1100$, hay dos cambios de signo.
\end{ejemplo}

\begin{figure}[ht]
\centering
\includegraphics[width=0.88\linewidth]{cap10/f_sorpresa}
\caption{Asíntotas contra curva real para $1000/[(s+1)(s^2+4s+100)]$: la resonancia empuja el cruce de ganancia hasta donde la fase ya pasó $-180\grados$.}
\label{fig:rf-sorpresa}
\end{figure}

\codigorepo{cap10_respuesta_frecuencial/Clase1_RespuestaFrecuencial/C1_Ej3_SegundoOrden.m}{sorpresa}{ejemplo \ref{ej:rf-sorpresa}}

\begin{profundiza}[El factor de segundo orden en detalle]
\textbf{Pico de resonancia.} El máximo de la magnitud y su frecuencia, para $\zeta<1/\sqrt2$, son
\begin{equation}
  \Mr=\frac{1}{2\zeta\sqrt{1-\zeta^2}},\qquad \wres=\wn\sqrt{1-2\zeta^2}.
  \label{eq:rf-Mr}
\end{equation}
Para $\zeta\geq1/\sqrt2\approx0.707$ no hay pico. La \cref{tab:rf-zeta} reúne los valores útiles, junto con el sobrepico de la respuesta al escalón, $\OS=100\,e^{-\pi\zeta/\sqrt{1-\zeta^2}}$.

\begin{center}\small
\begin{tabular}{@{}cccccc@{}}
\toprule
$\zeta$ & $20\log\frac{1}{2\zeta}$ (dB) & $\Mr$ & $\Mr$ (dB) & $\wres/\wn$ & $\OS$\\
\midrule
0.05 & 20.00 & 10.01 & 20.01 & 0.997 & 85.4\\
0.1  & 13.98 & 5.03 & 14.02 & 0.990 & 72.9\\
0.2  & 7.96  & 2.55 & 8.14  & 0.959 & 52.7\\
0.3  & 4.44  & 1.75 & 4.85  & 0.906 & 37.2\\
0.5  & 0.00  & 1.15 & 1.25  & 0.707 & 16.3\\
0.707& $-3.01$ & 1 & 0 & -- & 4.3\\
1    & $-6.02$ & 1 & 0 & -- & 0\\
\bottomrule
\end{tabular}
\captionof{table}{Corrección del factor de segundo orden en $\wn$, pico de resonancia y sobrepico al escalón.}
\label{tab:rf-zeta}
\end{center}

\textbf{La pendiente de la fase.} Cerca de $\wn$ la fase cae con pendiente
$\left.\frac{\dif\phi}{\dif\log_{10}\omega}\right|_{\wn}=-\frac{180\grados}{\pi}\,\frac{\ln10}{\zeta}\approx-\frac{132\grados}{\zeta}$ por década. Con $\zeta=0.2$ son $-660\grados$ por década: la fase cae casi como un escalón de $-180\grados$. Por eso una asíntota de fase de segundo orden dependería de $\zeta$, y por eso no la usamos.

\textbf{El factor de calidad.} En física y electrónica se usa $Q=1/(2\zeta)$, que es precisamente $|\Gma(\jw_n)|$. En un RLC serie $\zeta=\frac{R}{2}\sqrt{C/L}$, y en un sistema masa-resorte-amortiguador $\zeta=b/(2\sqrt{km})$.

\textbf{Leer $\zeta$ en un Bode medido.} La fase cruza $-90\grados$ en $\wn$; el pico da $\Mr$ y de ahí $\zeta$. Si la fase vale $-90\grados$ en $\omega=10$ y el pico es de $8.14\dB$, entonces $\wn=10$, $\Mr=10^{8.14/20}=2.55$ y $\zeta=0.20$ (para $\zeta$ pequeño, $\zeta\approx\frac{1}{2\Mr}=0.196$). La planta era $100/(s^2+4s+100)$.

\textbf{Ceros complejos: el filtro \emph{notch}.} Ceros sobre el eje imaginario anulan la ganancia en una frecuencia: $N(s)=\dfrac{s^2+\omega_0^2}{s^2+\frac{\omega_0}{Q}s+\omega_0^2}$ con $\omega_0=2\pi(60)=377$\,rad/s es el filtro que quita el zumbido de $60$\,Hz de la red en señales de ECG y de audio; $Q$ controla qué tan angosta es la muesca.
\end{profundiza}

\codigorepo{cap10_respuesta_frecuencial/Clase1_RespuestaFrecuencial/C1_Ej3_SegundoOrden.m}{leerzeta}{leer $\wn$ y $\zeta$ de un Bode medido}

%% ---------------------------------------------------------------------
\section{Ejemplo resuelto: $\Gma(s)=s/(s+10)$}
\label{sec:rf-ejemplo}
%% ---------------------------------------------------------------------

Este ejemplo recoge todo lo anterior y es el que más conviene dominar. Lo resolvemos en cinco pasos que sirven para cualquier malla abierta.

\begin{ejemplo}[label=ej:rf-ss10]{Diagrama asintótico de $s/(s+10)$}
Dibuje el diagrama de Bode asintótico de magnitud y fase de $\Gma(s)=s/(s+10)$ y compárelo con el real.
\solucion
\textbf{Paso 1: forma de Bode.} Se saca el $10$ del polo:
\[
  \Gma(s)=\frac{s}{10\big(\frac{s}{10}+1\big)}
  =\underbrace{\textcolor{ctrlazul}{\frac{1}{10}}}_{\text{ganancia}}\cdot\underbrace{\textcolor{ctrlcoral}{s}}_{\text{cero en el origen}}\cdot
  \underbrace{\textcolor{ctrlvioleta}{\frac{1}{\frac{s}{10}+1}}}_{\text{polo en }-10}.
\]
\textbf{Paso 2: identificar los factores.}
\begin{itemize}
  \item \textcolor{ctrlazul}{$1/10$}: $20\log\frac{1}{10}=-20\dB$ y $0\grados$ (recta horizontal).
  \item \textcolor{ctrlcoral}{$s$}: $+20\dBdec$ pasando por $0\dB$ en $\omega=1$, y $+90\grados$.
  \item \textcolor{ctrlvioleta}{$1/(s/10+1)$}: $0\dB$ hasta $\omega=10$ y $-20\dBdec$ después; fase de $0\grados$ a $-90\grados$ entre $\omega=1$ y $\omega=100$.
\end{itemize}
\textbf{Paso 3: dibujar cada factor por separado} (líneas delgadas de la \cref{fig:rf-ss10}).

\textbf{Paso 4: sumar.} En magnitud,
\[
  \omega<10:\ -20+20\log\omega,\qquad
  \omega>10:\ -20+20\log\omega-20\log\frac{\omega}{10}=0\dB .
\]
Desde $\omega=10$ las pendientes $+20$ y $-20$ se cancelan y la curva queda plana en $0\dB$. En $\omega=1$ la asíntota vale $-20\dB$ y en $\omega=0.1$, $-40\dB$. En fase, $\phi(\omega)=90\grados-\tan^{-1}(\omega/10)$: $90\grados$ para $\omega\le1$, $45\grados$ en $\omega=10$ y $0\grados$ para $\omega\ge100$.

\textbf{Paso 5: comparar con la curva real.} En la esquina,
\[
  \Big|\frac{j10}{10+j10}\Big|=\frac{1}{\sqrt2}\to-3\dB,\qquad\phi=90\grados-45\grados=45\grados.
\]
Otros valores reales: en $\omega=1$, $-20.04\dB$ y $84.3\grados$; en $\omega=20$, $-0.97\dB$ y $26.6\grados$; en $\omega=0.1$, $-40.0\dB$ y $89.4\grados$.
\end{ejemplo}

\begin{figure}[ht]
\centering
\includegraphics[width=\linewidth]{cap10/f_ej_s_s10}
\caption{Los tres factores de $s/(s+10)$ (delgados), su suma asintótica (negro) y la curva real (verde). Los puntos marcan la esquina: $-3\dB$ y $45\grados$.}
\label{fig:rf-ss10}
\end{figure}

\codigorepo{cap10_respuesta_frecuencial/Clase1_RespuestaFrecuencial/C1_Ej2_Asintotico_s_s10.m}{factores}{ejemplo \ref{ej:rf-ss10}}

\video{V10_asintotico_s_s10}{El diagrama se arma término a término con el mismo código de colores; al final, la curva real y los $-3\dB$ de la esquina.}

\paragraph{Interpretación: un filtro pasa-altos}
$s/(s+10)$ deja pasar con ganancia $1$ las frecuencias por encima de $10$\,rad/s y atenúa las de abajo ($-20\dB$ por cada década hacia la izquierda). En bajas frecuencias la salida, además de pequeña, va adelantada casi $90\grados$: el cero en el origen es un \emph{derivador}. Por ejemplo, para $x(t)=\sin(0.1t)$ la salida es $0.01\sin(0.1t+89.4\grados)$. Con una entrada compuesta $x(t)=\sin t+0.3\sin 20t$, la componente lenta sale multiplicada por $0.0995$ ($-20\dB$) y la rápida por $0.894$: en la salida domina la componente de $20$\,rad/s (\cref{fig:rf-pasaaltos}).

\begin{figure}[ht]
\centering
\includegraphics[width=\linewidth]{cap10/f_pasaaltos}
\caption{$s/(s+10)$ como filtro pasa-altos: de $\sin t+0.3\sin20t$ sobrevive casi solo la componente rápida.}
\label{fig:rf-pasaaltos}
\end{figure}

\video{V11_pasaaltos_s_s10}{Barrido de $0.1$ a $1000$\,rad/s mostrando entrada y salida junto al punto del Bode; al final, la entrada compuesta.}

\subsection{La receta y la hoja de trucos}

\begin{receta}[Diagrama de Bode asintótico, paso a paso]
\begin{enumerate}
  \item Llevar $\Gma(s)$ a la forma de Bode \eqref{eq:rf-formabode}: cada factor como $(s/\omega_i+1)$, y la ganancia $\Kb$ sola.
  \item Ordenar las esquinas de menor a mayor.
  \item Magnitud: arrancar con la recta de baja frecuencia, $20\log \Kb-20N\log\omega$, y recorrer de izquierda a derecha sumando $+20\dBdec$ por cada cero y $-20$ por cada polo ($-40$ por cada par complejo).
  \item Fase: factor por factor ($\pm45\grados$/dec entre $\omega_i/10$ y $10\,\omega_i$ para los reales; $-90N\grados$ de los integradores) y sumar.
  \item Verificar con \cmd{bode}. Con polos complejos y $\zeta$ pequeño, la curva real puede estar muy por encima de la asíntota.
\end{enumerate}
\end{receta}

\begin{anchocompleto}
\centering\small
\begin{tabular}{@{}llll@{}}
\toprule
Factor & Magnitud & Fase & Cuidado\\
\midrule
$\Kb$ & $20\log|\Kb|$, plana & $0\grados$ ($\pm180\grados$ si $\Kb<0$) & calcular $\Kb$ en forma de Bode\\
$1/s^N$ & $-20N\dBdec$, $0\dB$ en $\omega=1$ & $-90N\grados$ & define el tipo del sistema\\
$(s/\omega_1+1)^{\pm1}$ & $\pm20\dBdec$ desde $\omega_1$ & $0\to\pm90\grados$ ($\pm45\grados$ en $\omega_1$) & real: $\pm3\dB$ en $\omega_1$\\
polos complejos & $-40\dBdec$ desde $\wn$; pico $\approx\frac{1}{2\zeta}$ & $0\to-180\grados$ ($-90\grados$ en $\wn$) & curva real: \cmd{bode}\\
\bottomrule
\end{tabular}
\captionof{table}{Hoja de trucos de Bode.}
\label{tab:rf-trucos}
\end{anchocompleto}

\begin{pausa}
Sin dibujar: ¿cuánto vale la asíntota de magnitud de $\Gma(s)=\dfrac{100(s+1)}{s(s+10)}$ en $\omega=1$? ¿Y la curva real?
\respuesta{Forma de Bode: $\Gma=\dfrac{10(s+1)}{s(s/10+1)}$, $\Kb=10$. En $\omega=1$ la asíntota vale $20\log10-20\log1=20\dB$. La curva real está $3\dB$ por encima por el cero en $1$ (y un poco menos por el polo en $10$): $22.97\dB$.}
\end{pausa}

%% ---------------------------------------------------------------------
\section{Leer un Bode al revés: de la curva a la función de transferencia}
\label{sec:rf-identificacion}
%% ---------------------------------------------------------------------

La receta de la sección anterior también funciona en reversa. Si alguien mide el diagrama de magnitud de un sistema (con la máquina de Bode, por ejemplo), podemos reconstruir $\Gma(s)$ leyendo las pendientes y las esquinas de izquierda a derecha; la ganancia sale al final de un valor conocido.

\begin{ejemplo}[label=ej:rf-ident]{¿Qué función de transferencia es esta?}
El sistema es de fase mínima y de ganancia positiva. Su diagrama asintótico de magnitud (negro en la \cref{fig:rf-ident}) baja con $-20\dBdec$ en bajas frecuencias, se vuelve plano en $36\dB$ desde $\omega=1$ y cae con $-40\dBdec$ desde $\omega=10$. Determine $\Gma(s)$ y dibuje la fase.
\solucion Leemos las pendientes de izquierda a derecha:
\begin{itemize}
  \item $-20\dBdec$ desde el principio: un \textbf{polo en el origen}.
  \item En $\omega=1$ la pendiente pasa a $0$ ($-20+20$): un \textbf{cero en $-1$}.
  \item En $\omega=10$ la pendiente pasa a $-40$ ($0-40$): un \textbf{polo doble en $-10$}.
\end{itemize}
En forma de Bode, $\Gma(s)=\Kb\,\dfrac{s+1}{s\,(s/10+1)^2}$. Entre $\omega=1$ y $\omega=10$ la asíntota vale $20\log\Kb$ (el integrador y el cero se cancelan), de modo que
\[
  20\log \Kb=36\dB\ \Rightarrow\ \Kb=10^{36/20}=63.1,\qquad
  \Gma(s)=\frac{63.1\,(s+1)}{s\,(s/10+1)^2}=\frac{6309.6\,(s+1)}{s\,(s+10)^2}.
\]
La fase asintótica arranca en $-90\grados$ (integrador), sube hasta $-45\grados$ en $\omega=1$ por el cero, vuelve a $-90\grados$ en $\omega=10$ y termina en $-180\grados$ desde $\omega=100$.
\end{ejemplo}

\begin{figure}[ht]
\centering
\includegraphics[width=0.88\linewidth]{cap10/f_identificacion}
\caption{Identificación: el asintótico de magnitud es el dato; la fase sale ``gratis'' porque el sistema es de fase mínima.}
\label{fig:rf-ident}
\end{figure}

\codigorepo{cap10_respuesta_frecuencial/Clase2_Margenes/C2_Ej1_Identificacion.m}{identificacion}{ejemplo \ref{ej:rf-ident}}

\subsection{La fase sale gratis\ldots{} si el sistema es de fase mínima}
\label{sec:rf-faseminima}

En el ejemplo usamos una hipótesis fuerte: \emph{fase mínima}, es decir, sin ceros ni polos en el semiplano derecho. Bode demostró que, bajo esa hipótesis, la magnitud determina la fase~\citep{Bod40}. Sin ella, la magnitud no basta: $(1+s)/(1+2s)$ y $(1-s)/(1+2s)$ tienen \emph{exactamente} el mismo diagrama de magnitud, porque $|1+\jw|=|1-\jw|$, pero sus fases son muy distintas (\cref{fig:rf-nominima}). El cero en $+1$ agrega atraso de fase sin cambiar la magnitud.

\begin{figure}[ht]
\centering
\includegraphics[width=0.82\linewidth]{cap10/f_fase_no_minima}
\caption{Misma magnitud, distinta fase: el cero en el semiplano derecho (coral) atrasa la fase hasta $-180\grados$.}
\label{fig:rf-nominima}
\end{figure}

\begin{cuidado}
Antes de reconstruir la fase a partir de la magnitud, hay que saber (por la física del problema o por una medición de fase) que el sistema es de fase mínima. Los sistemas de fase no mínima aparecen en la práctica: la respuesta inversa de una caldera, un avión que cabecea, o el control de nivel con ciertas configuraciones. Para ellos el Bode de magnitud solo no basta y se requiere un estudio más detallado (criterio de Nyquist).
\end{cuidado}

\begin{pausa}
Una malla abierta de fase mínima tiene pendiente de $-40\dBdec$ en bajas frecuencias. ¿Cuántos integradores tiene y cuánto vale su fase a muy baja frecuencia?
\respuesta{Dos integradores ($-20\dBdec$ cada uno) y fase $-180\grados$. Es un sistema tipo 2 (\cref{sec:rf-tipo}).}
\end{pausa}

%% ---------------------------------------------------------------------
\section{Estabilidad desde la malla abierta: márgenes de ganancia y de fase}
\label{sec:rf-estabilidad}
%% ---------------------------------------------------------------------

\subsection{La frontera: $|\Gma|=1$ y fase $-180\grados$}

La ecuación característica del lazo cerrado es $1+\Gma(s)=0$, es decir, $\Gma(s)=-1$. Evaluada sobre el eje imaginario, la frontera entre estabilidad e inestabilidad es una frecuencia en la que, a la vez,
\begin{equation}
  |\Gma(\jw)|=1\ \ (0\dB)\qquad\text{y}\qquad \angle\Gma(\jw)=-180\grados .
  \label{eq:rf-frontera}
\end{equation}
La intuición es la de la \cref{sec:rf-lenguaje}: una senoide que da la vuelta al lazo regresa del mismo tamaño ($|\Gma|=1$) e invertida ($-180\grados$); el signo menos del comparador la vuelve a invertir, se suma a sí misma y el lazo oscila solo. Si en la frecuencia donde la fase vale $-180\grados$ la senoide regresa \emph{más pequeña} ($|\Gma|<1$), cada vuelta la achica y la oscilación muere.

\begin{ideaclave}
Si la malla abierta es estable y de fase mínima, el lazo cerrado es asintóticamente estable cuando la magnitud de $\Gma(\jw)$ es menor que $0\dB$ en la frecuencia en que su fase vale $-180\grados$.
\end{ideaclave}

\subsection{Márgenes de ganancia y de fase}

Saber que el lazo es estable no basta: queremos saber \emph{qué tan lejos} está de la frontera. Los dos márgenes miden esa distancia, cada uno en una dirección (\cref{fig:rf-margenes}).

\begin{definicion}[label=def:rf-gm]{Margen de ganancia}
El \clave{margen de ganancia} ($\GM$) es cuánto se puede aumentar la ganancia de la malla abierta, en decibeles, en la frecuencia donde la fase vale $-180\grados$ ($\wgm$, cruce de fase), antes de que el lazo cerrado se vuelva inestable:
\begin{equation}
  \GM=-20\log|\Gma(\jw_{GM})|\quad\text{dB},\qquad \angle\Gma(\jw_{GM})=-180\grados.
\end{equation}
\end{definicion}

\begin{definicion}[label=def:rf-pm]{Margen de fase}
El \clave{margen de fase} ($\PM$) es cuánto atraso de fase adicional se tolera en la frecuencia donde $|\Gma|=1$ ($\wpm$, cruce de ganancia) antes de llegar a $-180\grados$:
\begin{equation}
  \PM=180\grados+\angle\Gma(\jw_{PM}),\qquad |\Gma(\jw_{PM})|=1\ (0\dB).
\end{equation}
\end{definicion}

\begin{figure}[ht]
\centering
\includegraphics[width=0.88\linewidth]{cap10/f_margenes_def}
\caption{Lectura de los márgenes. El $\GM$ se mide en magnitud, en la frecuencia donde la fase cruza $-180\grados$; el $\PM$ se mide en fase, en la frecuencia donde la magnitud cruza $0\dB$.}
\label{fig:rf-margenes}
\end{figure}

\begin{cuidado}[Los nombres de MATLAB invitan a confundirse]
En \cmd{[Gm,Pm,Wcg,Wcp] = margin(G)}, \cmd{Wcg} (\emph{gain crossover} en otros textos) es en realidad la frecuencia donde se mide el $\GM$, es decir, nuestro $\wgm$, y \cmd{Wcp} es la frecuencia del $\PM$, nuestro $\wpm$. Además \cmd{Gm} sale en \emph{veces}: el $\GM$ en dB es \cmd{20*log10(Gm)}. La función del curso \cmd{margenesTabla(G)} imprime todo ya convertido.
\end{cuidado}

\begin{ejemplo}[label=ej:rf-margenes]{Márgenes asintóticos contra exactos}
Determine el $\GM$ y el $\PM$ de $\Gma(s)=\dfrac{10\,000}{(s+1)(s+10)(s+100)}$ con el diagrama asintótico y compare con \cmd{margin}.
\solucion \textbf{Forma de Bode:} $\Gma(s)=\dfrac{10}{(s+1)(s/10+1)(s/100+1)}$, es decir, $20\dB$ en bajas frecuencias y tres esquinas.

\textbf{Cruce por $0\dB$ (PM).} La asíntota baja $-20\dBdec$ desde $\omega=1$ y llega a $0\dB$ justo en $\omega=10$. La fase asintótica allí: el polo en $1$ ya aportó $-90\grados$, el de $10$ está en su esquina ($-45\grados$) y el de $100$ todavía no empieza. Total $-135\grados$: $\PM\approx45\grados$.

\textbf{Cruce por $-180\grados$ (GM).} Entre $\omega=10$ y $\omega=100$ la fase asintótica cae $-45\grados$ por década por el polo en $10$ y otros $-45\grados$ por década por el de $100$, desde $-135\grados$: llega a $-180\grados$ media década después de $10$, en $\wgm=10^{1.5}=31.6$\,rad/s. La magnitud asintótica allí es $0-40\times0.5=-20\dB$: $\GM\approx20\dB$.

\textbf{Comparación.}
\begin{center}\small
\begin{tabular}{@{}lcc@{}}
\toprule
 & asintótico & exacto (\cmd{margin})\\
\midrule
$\PM$ & $45\grados$ ($\omega=10$) & $54.9\grados$ ($\omega=7.80$)\\
$\GM$ & $20\dB$ ($\omega=31.6$) & $21.7\dB$ ($\omega=33.3$)\\
\bottomrule
\end{tabular}
\end{center}
El asintótico es conservador aquí: sirve para estimar en el tablero; el número final siempre con \cmd{margin}.
\end{ejemplo}

\begin{figure}[ht]
\centering
\includegraphics[width=0.9\linewidth]{cap10/f_margenes_asint}
\caption{Ejemplo~\ref{ej:rf-margenes}: los cuatro factores, su suma asintótica y la curva real, con los márgenes asintóticos marcados.}
\label{fig:rf-margenes-asint}
\end{figure}

\codigorepo{cap10_respuesta_frecuencial/Clase2_Margenes/C2_Ej2_Margenes.m}{margenes}{ejemplo \ref{ej:rf-margenes}}
\video{V3_asintotico_margenes}{El Bode se arma término a término; al final, PM y GM sobre la asintótica.}

%% ---------------------------------------------------------------------
\subsection{¿Qué significan los márgenes en la práctica?}
\label{sec:rf-significado}

\paragraph{El GM: el micrófono y el parlante}
\notamargen{El ``aullido'' de un sistema de sonido es un lazo cerrado que perdió su margen de ganancia.}
Un micrófono frente a un parlante forma un lazo: lo que sale del parlante vuelve a entrar al micrófono. Subir el volumen es subir la ganancia de la malla abierta \emph{sin cambiar la fase}. Mientras haya margen de ganancia no pasa nada; cuando el volumen sube más que el $\GM$, en la frecuencia donde la fase acústica vale $-180\grados$ aparece el aullido: una oscilación sostenida. El $\GM$ es, literalmente, cuánto se puede subir el volumen.

\paragraph{El PM: un presupuesto de fase para lo que no modelamos}
Ningún modelo es perfecto. Sensores, actuadores y filtros que dejamos por fuera agregan atraso de fase que la malla abierta modelada no ve. El $\PM$ es el \emph{presupuesto} con el que pagamos ese atraso. Para verlo, pongamos un sensor de primer orden $\Hs(s)=a/(s+a)$ en la retroalimentación del ejemplo~\ref{ej:rf-margenes} (\cref{fig:rf-lazosensor}). En el cruce, el sensor resta aproximadamente $\tan^{-1}(\wpm/a)$, de modo que
\begin{equation}
  \PM_{\text{nuevo}}\approx\PM-\tan^{-1}\frac{\wpm}{a}\qquad(\text{buena estimación si } a\gg\wpm).
\end{equation}

\begin{figure}[ht]
\centering
\includegraphics[width=0.55\linewidth]{cap10/lazo_sensor}
\caption{El lazo del ejemplo~\ref{ej:rf-margenes} con un sensor de primer orden. La malla abierta es ahora $\Gma=G\,\Hs$.}
\label{fig:rf-lazosensor}
\end{figure}

\begin{center}\small
\begin{tabular}{@{}ccccc@{}}
\toprule
$a$ (rad/s) & $\PM$ estimado & $\PM$ (\cmd{margin}) & $\GM$ & sobrepico\\
\midrule
sin sensor & -- & $54.9\grados$ & $21.7\dB$ & $18\%$\\
$50$ & $46.0\grados$ & $46.5\grados$ & $12.7\dB$ & $28\%$\\
$20$ & $33.6\grados$ & $36.4\grados$ & $8.6\dB$ & $44\%$\\
$10$ & $17.0\grados$ & $26.3\grados$ & $6.1\dB$ & $68\%$\\
$5$  & $-2.4\grados$ & $18.7\grados$ & $4.7\dB$ & $104\%$\\
\bottomrule
\end{tabular}
\captionof{table}{El sensor ``gasta'' fase. Con $a=50$, casi diez veces más rápido que el cruce, la estimación es buena; con sensores lentos la estimación es pesimista (el sensor también baja la magnitud y corre el cruce a la izquierda), pero la tendencia es la misma.}
\label{tab:rf-sensor}
\end{center}

\begin{figure}[ht]
\centering
\includegraphics[width=\linewidth]{cap10/f_sensor}
\caption{Un sensor más lento gasta más fase en el cruce (derecha) y el lazo cerrado se vuelve más oscilatorio (izquierda), aunque la planta sea la misma.}
\label{fig:rf-sensor}
\end{figure}

\codigorepo{cap10_respuesta_frecuencial/Clase2_Margenes/C2_Video_SensorMargenFase.m}{sensor}{tabla \ref{tab:rf-sensor}}
\video{V12_sensor_margen_fase}{El polo del sensor baja de $300$ a $5$\,rad/s: se ve cómo se consume el PM y crece el sobrepico.}

\begin{ideaclave}
El $\GM$ dice cuánta ganancia de más aguanta el lazo; el $\PM$, cuánto atraso de fase no modelado aguanta. Un buen diseño necesita los dos.
\end{ideaclave}

\begin{profundiza}[Los márgenes como distancia al punto $-1$]
Ambos márgenes miden qué tan cerca pasa $\Gma(\jw)$ del punto crítico $-1$, cada uno en una dirección: el $\GM$ sobre el eje real negativo y el $\PM$ sobre el círculo unitario. Una medida que los resume es la distancia mínima $\min_\omega|1+\Gma(\jw)|=1/M_s$, donde $M_s=\max_\omega|\Sens(\jw)|$ es el pico de la sensitividad $\Sens=1/(1+\Gma)$. Un $M_s$ entre $1.3$ y $2$ garantiza a la vez $\GM\geq20\log\frac{M_s}{M_s-1}$ y $\PM\geq2\sin^{-1}\frac{1}{2M_s}$; por ejemplo, $M_s=2$ asegura $\GM\geq6\dB$ y $\PM\geq29\grados$~\citep{AsMu10}.
\end{profundiza}

%% ---------------------------------------------------------------------
\subsection{¿Cuánto margen es suficiente?}

\begin{itemize}
  \item \textbf{Fase mínima:} $\GM$ y $\PM$ positivos $\Rightarrow$ estable. Márgenes negativos $\Rightarrow$ inestable.
  \item \textbf{Regla práctica:} $30\grados<\PM<60\grados$ y $\GM>6\dB$ dan una estabilidad que tolera variaciones razonables de la ganancia y de las constantes de tiempo.
  \item \textbf{Fase no mínima} (ceros o polos en el semiplano derecho): el Bode solo no basta; se requiere un estudio más detallado (criterio de Nyquist).
\end{itemize}

\begin{pausa}
Tres lazos de fase mínima: (A) $\PM=40\grados$, $\GM=12\dB$; (B) $\PM=-15\grados$, $\GM=-4\dB$; (C) $\PM=70\grados$, $\GM=2\dB$. ¿Cuáles son estables? ¿Cuál le preocupa más?
\respuesta{(A) estable, con márgenes cómodos. (B) inestable: hay que bajar la ganancia al menos $4\dB$. (C) estable pero frágil: $2\dB$ es un factor $10^{2/20}=1.26$, así que un $26\%$ más de ganancia lo inestabiliza. Un $\PM$ grande no basta: hay que mirar los dos márgenes.}
\end{pausa}

%% ---------------------------------------------------------------------
\subsection{¿Cuánta ganancia aguanta? Routh-Hurwitz contra Bode}
\label{sec:rf-rangokc}

\begin{figure}[ht]
\centering
\includegraphics[width=0.62\linewidth]{cap10/lazo_basico}
\caption{Lazo de control con retroalimentación. En los ejemplos de diseño, $\Hs(s)=1$.}
\label{fig:rf-lazo}
\end{figure}

\begin{ejemplo}[label=ej:rf-rangokc]{Rango de $\Kc$ por dos caminos}
En el lazo de la \cref{fig:rf-lazo}, con $\Hs=1$, $\Gc(s)=\Kc$ y $\Gp(s)=\dfrac{1}{(s+1)(s+10)(s+100)}$, determine el rango de $\Kc$ que garantiza estabilidad asintótica con Routh-Hurwitz y con el diagrama de Bode.
\solucion \textbf{Routh-Hurwitz.} El polinomio característico es $s^3+111s^2+1110s+(1000+\Kc)$. La fila $s^1$ exige $111\cdot1110-(1000+\Kc)>0$ y la fila $s^0$, $1000+\Kc>0$:
\[
  -1000<\Kc<122\,210 .
\]
En $\Kc=\Kcr=122\,210$ el polinomio auxiliar $111s^2+123\,210=0$ da $s=\pm j\sqrt{1110}$: el lazo oscila con $\wcr=33.32$\,rad/s y periodo $\Pcr=2\pi/\wcr=0.189$\,s.

\textbf{Bode.} Para dibujar hay que fijar algún $\Kc$. Conviene el que deja la ganancia DC en $0\dB$, $\Kc=1000$, con lo que $\Gma(s)=\dfrac{1}{(s+1)(s/10+1)(s/100+1)}$. El cruce de fase asintótico está otra vez en $\omega=31.6$, donde la magnitud asintótica vale $-40\dB$: hay $\GM\approx+40\dB$, es decir, la ganancia se puede multiplicar por $100$. Así $\Kcr\approx1000\times100=100\,000$ y el rango positivo es $0<\Kc<100\,000$. Con \cmd{margin} el $\GM$ exacto es $41.74\dB$ ($\times122.2$), que reproduce $\Kcr=122\,210$.

La lectura asintótica se equivoca hacia el lado seguro: subestima la ganancia crítica.
\end{ejemplo}

\begin{figure}[ht]
\centering
\includegraphics[width=0.9\linewidth]{cap10/f_rango_kc}
\caption{Respuesta al escalón del lazo cerrado al subir $\Kc$: en $\Kc=\Kcr$ la oscilación ya no se amortigua.}
\label{fig:rf-rangokc}
\end{figure}

\codigorepo{cap10_respuesta_frecuencial/Clase2_Margenes/C2_Ej3_RangoKc.m}{routh}{Routh-Hurwitz sin Symbolic Toolbox}
\video{V4_barrido_Kc}{Subir $\Kc$ solo sube la magnitud: la barra roja del GM se acorta hasta cero en $\Kc=122\,210$.}

\begin{modelo}[LazoKc.slx]
\cmd{crearModelo\_LazoKc.m} arma el lazo escalón $\to$ suma $\to$ $\Kc$ $\to$ $\Gp(s)$ con retroalimentación unitaria. Pruebe $\Kc=1000$, $50\,000$, $122\,210$ (marginal) y $150\,000$ (inestable). Para otra planta cambie el parámetro \emph{Poles} del bloque $\Gp$.
\end{modelo}

\paragraph{Conexión con lo que ya vimos: Ziegler-Nichols}
\notamargen{Ziegler y Nichols publicaron sus reglas en 1942~\citep{ZiNi42}.}
El $\GM$ entrega la ganancia crítica y $\wgm$ la frecuencia crítica: son justo los datos del método de oscilación sostenida de Ziegler y Nichols (capítulo de sintonización). Con nuestro ejemplo, $\Kcr=122\,210$ y $\Pcr=0.189$\,s, y la tabla de ZN para un PID da
\[
  \Kc=0.6\,\Kcr=73\,326,\qquad \Ti=\Pcr/2=0.0943\ \text{s},\qquad \Td=\Pcr/8=0.0236\ \text{s}.
\]
La ventaja del Bode es que \emph{no hay que llevar la planta real al borde de la inestabilidad} para medir $\Kcr$ y $\Pcr$: se leen en el diagrama. En la \cref{sec:rf-pid} veremos qué margen de fase deja esa sintonización.

%% ---------------------------------------------------------------------
\section{Del Bode de malla abierta al transitorio del lazo cerrado}
\label{sec:rf-lazocerrado}
%% ---------------------------------------------------------------------

Los márgenes dicen \emph{si} el lazo cerrado es estable. Ahora queremos saber \emph{cómo} se comporta: cuánto sobrepasa, qué tan rápido responde. Lo haremos con el caso de referencia que conocemos desde el capítulo de segundo orden y luego lo usaremos como aproximación de polos dominantes.

\subsection{El caso de referencia}
Con $\Gc=1$ y $\Gp(s)=\dfrac{\wn^2}{s(s+2\zeta\wn)}$, la malla abierta es tipo 1 y el lazo cerrado es el segundo orden estándar,
\begin{equation}
  \Tlc(s)=\frac{Y(s)}{R(s)}=\frac{\wn^2}{s^2+2\zeta\wn s+\wn^2}.
  \label{eq:rf-T2}
\end{equation}

\subsection{El lazo cerrado visto en frecuencia: pico de resonancia y ancho de banda}

El diagrama de magnitud del \emph{lazo cerrado} $|\Tlc(\jw)|$ tiene dos rasgos que resumen su comportamiento (\cref{fig:rf-Tzeta}).

\paragraph{Pico de resonancia}
Maximizar $|\Tlc(\jw)|=\wn^2/\sqrt{(\wn^2-\omega^2)^2+4\zeta^2\wn^2\omega^2}$ equivale a minimizar el radicando como función de $\omega^2$; derivando e igualando a cero se obtienen el pico $\Mr$ y su frecuencia $\wres$ de la ecuación~\eqref{eq:rf-Mr}:
\[
  \Mr=\frac{1}{2\zeta\sqrt{1-\zeta^2}},\qquad \wres=\wn\sqrt{1-2\zeta^2},\qquad 0<\zeta<\tfrac{1}{\sqrt2}.
\]
Un $\Mr$ grande indica polos dominantes poco amortiguados y una respuesta oscilatoria; para $\zeta>0.707$ no hay pico ($\Mr=1$).

\begin{cuidado}
No confundir el pico de resonancia $\Mr$ (frecuencia, lazo cerrado) con el sobrepico $\OS$ (tiempo, respuesta al escalón). Están relacionados porque ambos dependen solo de $\zeta$, pero son números distintos: con $\zeta=0.5$, $\Mr=1.15$ ($1.25\dB$) y $\OS=16.3\%$. En versiones anteriores de estas notas el pico de resonancia se llamaba $M_P$; para evitar la confusión con el sobrepico, ahora es $\Mr$.
\end{cuidado}

\paragraph{Ancho de banda}
El \clave{ancho de banda} $\wbw$ es la frecuencia donde $|\Tlc|$ cae $3\dB$ respecto a su valor de baja frecuencia, $|\Tlc(\jw_{BW})|=1/\sqrt2$. Mide la habilidad del lazo para reproducir fielmente la referencia:
\begin{equation}
  \wbw=\wn\sqrt{(1-2\zeta^2)+\sqrt{4\zeta^4-4\zeta^2+2}} .
  \label{eq:rf-BW}
\end{equation}
Como $\ts=4/(\zeta\wn)$ y $\tp=\pi/(\wn\sqrt{1-\zeta^2})$, para un $\zeta$ dado un ancho de banda mayor significa un sistema más rápido. Más ancho de banda también deja pasar más ruido de medición: nada es gratis.

\begin{figure}[ht]
\centering
\includegraphics[width=0.85\linewidth]{cap10/f_lazo_cerrado_zeta}
\caption{Magnitud del lazo cerrado $|\Tlc(\jw)|$ para varios $\zeta$: el pico $\Mr$ crece al bajar $\zeta$; la línea punteada marca los $-3\dB$ que definen $\wbw$.}
\label{fig:rf-Tzeta}
\end{figure}

\subsection{Del margen de fase a $\zeta$, paso a paso}

Ahora la idea central de la sección: obtener $\zeta$ \emph{sin} cerrar el lazo, mirando solo la malla abierta $\Gma(s)=\wn^2/[s(s+2\zeta\wn)]$.

\textbf{Paso 1: frecuencia de cruce.} Imponiendo $|\Gma(\jw_1)|=\wn^2/|{-\omega_1^2}+j2\zeta\wn\omega_1|=1$ y despejando,
\begin{equation}
  \omega_1=\wn\sqrt{-2\zeta^2+\sqrt{1+4\zeta^4}} .
  \label{eq:rf-omega1}
\end{equation}
\textbf{Paso 2: fase en el cruce y margen de fase.} $\angle\Gma(\jw_1)=-90\grados-\tan^{-1}\dfrac{\omega_1}{2\zeta\wn}$, de donde
\begin{equation}
  \PM=180\grados+\angle\Gma(\jw_1)=\tan^{-1}\frac{2\zeta}{\sqrt{-2\zeta^2+\sqrt{1+4\zeta^4}}}.
  \label{eq:rf-PMzeta}
\end{equation}
El margen de fase depende \emph{solo} de $\zeta$. La \cref{fig:rf-pmzeta} es la curva que usaremos todo el capítulo; para $\zeta<0.6$ vale la aproximación $\PM\approx100\,\zeta$.

\begin{figure}[ht]
\centering
\includegraphics[width=0.85\linewidth]{cap10/f_pm_zeta}
\caption{Margen de fase y sobrepico en función de $\zeta$.}
\label{fig:rf-pmzeta}
\end{figure}

\begin{center}\small
\begin{tabular}{@{}cccccc@{}}
\toprule
$\zeta$ & 0.25 & 0.40 & 0.50 & 0.60 & 0.707\\
\midrule
$\PM$ & $28.0\grados$ & $43.1\grados$ & $51.8\grados$ & $59.2\grados$ & $65.5\grados$\\
$\OS$ & $44.4\%$ & $25.4\%$ & $16.3\%$ & $9.5\%$ & $4.3\%$\\
\bottomrule
\end{tabular}
\captionof{table}{Valores de referencia de la relación $\PM\leftrightarrow\zeta\leftrightarrow\OS$. Un $\PM\geq65.5\grados$ asegura que el lazo cerrado no tiene pico de resonancia.}
\label{tab:rf-pmzeta}
\end{center}

\video{V5_PM_zeta}{Al subir $K$ en $K/[s(s+10)]$ el cruce se va a la derecha, el PM baja y el punto rojo recorre la curva $\PM(\zeta)$ mientras el escalón se vuelve más oscilatorio.}

\begin{receta}[Del Bode de malla abierta al transitorio]
\begin{enumerate}
  \item \textbf{Leer el PM} y su frecuencia $\wpm=\omega_1$ en la malla abierta.
  \item \textbf{Obtener $\zeta$} de \eqref{eq:rf-PMzeta} (función \cmd{zetaDePM}, o $\zeta\approx\PM/100$).
  \item \textbf{Obtener $\wn$} de \eqref{eq:rf-omega1}.
  \item \textbf{Transitorio:} $\OS=100\,e^{-\pi\zeta/\sqrt{1-\zeta^2}}$, $\tp=\pi/(\wn\sqrt{1-\zeta^2})$, $\ts=4/(\zeta\wn)$; o $\wbw$ con \eqref{eq:rf-BW}.
\end{enumerate}
Es exacto para el segundo orden de referencia; para órdenes mayores es una aproximación de polos dominantes. Siempre verificar con \cmd{step}.
\end{receta}

\begin{ejemplo}[label=ej:rf-pmzeta]{$\OS$, $\tp$ y $\ts$ sin simular}
Para $\Gma(s)=\dfrac{100}{s(s+10)}$ con retroalimentación unitaria, determine $\OS$, $\tp$ y $\ts$ del lazo cerrado a partir de la malla abierta.
\solucion Con \cmd{margin}: $\PM=51.8\grados$ en $\wpm=7.86$\,rad/s. De \eqref{eq:rf-PMzeta}, $\zeta=0.50$; de \eqref{eq:rf-omega1}, $\wn=7.86/0.786=10$\,rad/s. Entonces
\[
  \OS=16.3\%,\qquad \tp=\frac{\pi}{10\sqrt{1-0.25}}=0.363\ \text{s},\qquad \ts=\frac{4}{0.5\cdot10}=0.8\ \text{s},
\]
y además $\Mr=1.25\dB$ en $\wres=7.07$\,rad/s y $\wbw=12.7$\,rad/s. Como este sistema \emph{es} el segundo orden de referencia ($\wn^2=100$, $2\zeta\wn=10$), los valores son exactos (\cref{fig:rf-ejpmzeta}).

\textbf{¿Y si solo tengo el asintótico?} La asíntota $10/[s(s/10+1)]$ cruza $0\dB$ en $\omega=10$, donde la fase asintótica es $-135\grados$: $\PM\approx45\grados$, $\zeta\approx0.42$ y $\OS\approx23\%$. Son unos siete puntos de error en el sobrepico. La razón: cerca de la esquina $\omega=10$ la magnitud real está $3\dB$ por debajo de la asíntota, así que el cruce real ocurre antes ($7.86$) y con más fase.
\end{ejemplo}

\begin{figure}[ht]
\centering
\includegraphics[width=\linewidth]{cap10/f_ej_pm_zeta}
\caption{Ejemplo~\ref{ej:rf-pmzeta}: el PM de la malla abierta predice el sobrepico del lazo cerrado.}
\label{fig:rf-ejpmzeta}
\end{figure}

\codigorepo{cap10_respuesta_frecuencial/Clase3_LazoCerrado/C3_Ej1_PM_zeta.m}{receta}{ejemplo \ref{ej:rf-pmzeta}}

\begin{pausa}
Repita el ejemplo para $\Gma(s)=100/[s(s+5)]$. ¿Qué cambió y por qué?
\respuesta{Ahora $\wn=10$ y $2\zeta\wn=5$, así que $\zeta=0.25$: $\PM=28.0\grados$, $\OS=44.4\%$, $\ts=4/(0.25\cdot10)=1.6$\,s y $\tp=0.32$\,s. El polo de la planta está más cerca del origen: en el cruce aporta más atraso de fase, el PM baja y el lazo cerrado oscila más.}
\end{pausa}

\begin{ideaclave}
\begin{itemize}
  \item El $\OS$ baja al aumentar el $\PM$; la velocidad sube al aumentar $\wbw$ (en la práctica, al correr el cruce $\wpm$ a la derecha).
  \item En diseño es más común especificar $\PM$ y $\GM$ que $\Mr$.
  \item El error en estado estacionario mejora al subir la magnitud en \emph{bajas} frecuencias; el transitorio depende de lo que pasa \emph{en el cruce}. Esa separación es la que explotan los compensadores (\cref{sec:rf-compensadores}).
\end{itemize}
\end{ideaclave}

%% ---------------------------------------------------------------------
\section{Tipo y error: todo está en las bajas frecuencias}
\label{sec:rf-tipo}
%% ---------------------------------------------------------------------

Con retroalimentación unitaria, el error de seguimiento $E(s)=R(s)-Y(s)$ en estado estacionario depende del número de integradores de la malla abierta (el \emph{tipo}) y de una constante estática (\cref{cap:er}). Las dos cosas se leen en el extremo izquierdo del diagrama de magnitud (\cref{fig:rf-tipos}):

\begin{itemize}
  \item \textbf{Tipo 0:} la magnitud arranca plana. El nivel de baja frecuencia es $20\log\kp$.
  \item \textbf{Tipo 1:} arranca con $-20\dBdec$. La asíntota de baja frecuencia, $\kv/\omega$, cruza $0\dB$ en $\omega=\kv$.
  \item \textbf{Tipo 2:} arranca con $-40\dBdec$. La asíntota, $\ka/\omega^2$, cruza $0\dB$ en $\omega=\sqrt{\ka}$.
\end{itemize}

\begin{figure}[ht]
\centering
\includegraphics[width=\linewidth]{cap10/f_tipos}
\caption{El tipo se ve en la pendiente inicial y la constante estática, en un nivel o en un cruce de la asíntota de baja frecuencia.}
\label{fig:rf-tipos}
\end{figure}

\begin{ejemplo}[label=ej:rf-tipos]{Lea las constantes y el error}
Los tres diagramas de la \cref{fig:rf-tipos} corresponden a $\dfrac{0.01}{s/10+1}$, $\dfrac{1}{s(s/10+1)}$ y $\dfrac{100}{s^2(s/100+1)}$. Determine el tipo, la constante estática y el error ante escalón, rampa y parábola unitarios (los tres lazos cerrados son estables).
\solucion
\begin{center}\small
\begin{tabular}{@{}lccccc@{}}
\toprule
Sistema & lectura en el Bode & constante & escalón & rampa & parábola\\
\midrule
Tipo 0 & nivel $-40\dB$ & $\kp=10^{-40/20}=0.01$ & $\frac{1}{1+\kp}=0.99$ & $\infty$ & $\infty$\\
Tipo 1 & cruce en $\omega=1$ & $\kv=1$ & $0$ & $\frac{1}{\kv}=1$ & $\infty$\\
Tipo 2 & cruce en $\omega=10$ & $\ka=10^2=100$ & $0$ & $0$ & $\frac{1}{\ka}=0.01$\\
\bottomrule
\end{tabular}
\end{center}
\end{ejemplo}

\codigorepo{cap10_respuesta_frecuencial/Clase3_LazoCerrado/C3_Ej2_TipoSistema.m}{tipos}{ejemplo \ref{ej:rf-tipos}}

\begin{ideaclave}
Para mejorar el error hay que \emph{subir la curva en bajas frecuencias}; para el transitorio importa la \emph{zona del cruce}. Un compensador en atraso explota exactamente esta separación (\cref{sec:rf-atraso}).
\end{ideaclave}

\begin{pausa}
¿Por qué en un sistema tipo 2 la asíntota de baja frecuencia cruza $0\dB$ en $\sqrt{\ka}$ y no en $\ka$?
\respuesta{En bajas frecuencias $\Gma(\jw)\approx\ka/(\jw)^2$, cuya magnitud es $\ka/\omega^2$. Igualándola a $1$ resulta $\omega^2=\ka$, es decir, $\omega=\sqrt{\ka}$.}
\end{pausa}

%% ---------------------------------------------------------------------
\section{Diseño ajustando la ganancia}
\label{sec:rf-ajuste}
%% ---------------------------------------------------------------------

Primer diseño, con una sola perilla. Como una ganancia sube o baja la curva de magnitud \emph{sin tocar la fase} (factor 1), podemos escoger la frecuencia donde la fase ya tiene el margen que queremos y luego subir o bajar la magnitud hasta que esa frecuencia sea el cruce.

\begin{receta}[Ajuste de ganancia para un $\OS$ dado]
\begin{enumerate}
  \item De la especificación: $\OS\Rightarrow\zeta=\dfrac{-\ln(\OS/100)}{\sqrt{\pi^2+\ln^2(\OS/100)}}\Rightarrow\PMdes$ con \eqref{eq:rf-PMzeta}.
  \item Dibujar el Bode con una ganancia cómoda y buscar la frecuencia $\omega_d$ donde $\angle\Gma(\jw_d)=-180\grados+\PMdes$.
  \item Subir o bajar la magnitud hasta que valga $0\dB$ en $\omega_d$: $\Kc=1/|\Gp(\jw_d)|$.
  \item Verificar con \cmd{margin} y \cmd{step}.
\end{enumerate}
\end{receta}

\begin{ejemplo}[label=ej:rf-ajuste]{$\Kc$ para un $20\%$ de sobrepico}
En el lazo de la \cref{fig:rf-lazo} con $\Gc=\Kc$ y $\Gp(s)=\dfrac{1}{s(s+20)(s+100)}$, determine $\Kc$ para que el lazo cerrado tenga $\OS=20\%$.
\solucion Es un sistema de tercer orden y tipo 1. En lugar de despreciar un polo para parecerse a un segundo orden, usamos la receta directamente.

\textbf{Paso 1.} $\zeta=0.456$ y $\PMdes=48.15\grados$: la fase objetivo es $-180\grados+48.15\grados=-131.85\grados$.

\textbf{Paso 2.} Con $\Kc=2000$ (ganancia DC de los polos reales normalizada) se busca la frecuencia con fase $-131.85\grados$: $\omega_d=13.56$\,rad/s. Allí la magnitud vale $-24.36\dB$.

\textbf{Paso 3.} Hay que subir $24.36\dB$: $\Kc=2000\times10^{24.364/20}\approx33\,056$.

\textbf{Paso 4.} Con este $\Kc$ el lazo cerrado tiene $\OS=19.9\%$ y $\tp=0.21$\,s (\cref{fig:rf-ajuste}). El sistema es de tercer orden, pero el polo en $-100$ está lejos: domina el par complejo.
\end{ejemplo}

\begin{figure}[ht]
\centering
\includegraphics[width=\linewidth]{cap10/f_ajuste_ganancia}
\caption{Ejemplo~\ref{ej:rf-ajuste}: se busca la frecuencia con la fase objetivo y se sube la curva $24.36\dB$ para que allí esté el cruce.}
\label{fig:rf-ajuste}
\end{figure}

\codigorepo{cap10_respuesta_frecuencial/Clase3_LazoCerrado/C3_Ej3_AjusteGanancia.m}{ajuste}{ejemplo \ref{ej:rf-ajuste}: sin buscar a ojo}
\video{V6_ajuste_ganancia}{La curva de magnitud sube hasta cruzar $0\dB$ en $13.56$\,rad/s mientras la fase queda quieta.}

\begin{cuidado}[Leer a ojo cuesta cifras]
Leyendo el diagrama a mano se obtiene algo como $-24.2\dB$ y $\Kc=2000\times10^{24.2/20}\approx32\,400$ (el handout anterior reportaba $33\,192$ para esa lectura: era un error aritmético); el valor exacto, buscando la frecuencia con \cmd{fzero}, es $\Kc=33\,056$. La diferencia no importa para el diseño, pero sí para reportar: diga siempre de dónde sale cada cifra.
\end{cuidado}

\subsection{El límite del ajuste de ganancia}

Con una sola perilla solo se cumple \emph{una} especificación. Con $\Kc=33\,056$ la constante de velocidad es $\kv=\Kc/2000=16.5$. Si además se exigiera $\kv=40$, haría falta $\Kc=80\,000$, y entonces $\PM=25.4\grados$ y $\OS\approx49\%$: subir $\Kc$ para el error destruye el margen de fase. Necesitamos cambiar la \emph{forma} de la curva, no solo subirla o bajarla. Eso es un compensador.

%% ---------------------------------------------------------------------
\section{Compensadores en adelanto, en atraso y atraso-adelanto}
\label{sec:rf-compensadores}
%% ---------------------------------------------------------------------

\subsection{El dilema y la filosofía de diseño}

En el lazo de la \cref{fig:rf-lazo} con $\Gp(s)=K/[s(s+2)]$ queremos dos cosas: un error del $5\%$ ante una rampa ($\kv=20$, es decir, $K=40$) y un margen de fase de $45\grados$. Con $K=40$ el margen de fase es apenas $18\grados$ (sobrepico del $60\%$); bajando $K$ se recupera el margen pero se pierde el error. Con una sola perilla no alcanza: necesitamos cambiar la forma de la curva.

\begin{historia}[De las redes RC a los servomecanismos]
En Bell Labs las redes de adelanto y atraso eran simples combinaciones de resistencias y condensadores dentro del lazo de los amplificadores. Durante la Segunda Guerra Mundial los servomecanismos de radares y cañones se diseñaron en frecuencia en el MIT Radiation Laboratory, de donde salió el libro de James, Nichols y Phillips~\citep{JaNiPh47}. Hoy las mismas redes se implementan con amplificadores operacionales o en software: un adelanto es ``un PD con filtro'' y un atraso es ``casi un PI''. Lo que no ha cambiado es que se diseña mirando el Bode de $\Gc\Gp$: \emph{baja frecuencia = error; cruce = transitorio y estabilidad; alta frecuencia = ruido}.
\end{historia}

\begin{receta}[Cómo se reparte el trabajo del compensador]
\begin{enumerate}
  \item \textbf{$\Kc$ para el error.} Se escoge la ganancia DC más pequeña que cumpla $\ess$ (con $\kp$, $\kv$ o $\ka$).
  \item \textbf{$\Gcp(s)$ para la forma.} Se agregan términos con ganancia DC unitaria, $\Gcp(0)=1$: no cambian el error, pero sí el $\PM$ y la velocidad. El compensador es $\Gc(s)=\Kc\,\Gcp(s)$.
  \item \textbf{Verificar} en el Bode de $\Gc\Gp$ ($\PM$, $\GM$) y en el tiempo (escalón, rampa, esfuerzo de control).
\end{enumerate}
Supuestos: $\Gc\Gp$ sin polos en el semiplano derecho ni cancelaciones en el eje imaginario; así basta el Bode para decidir la estabilidad.
\end{receta}

%% ---------------------------------------------------------------------
\subsection{El compensador en adelanto: agregar fase justo donde hace falta}
\label{sec:rf-adelanto}

\begin{equation}
  \Gcp(s)=\frac{1+s/\wL}{1+s/\wH}=\alpha\,\frac{s+\wL}{s+\wH},\qquad
  \alpha=\frac{\wH}{\wL}>1 .
  \label{eq:rf-adelanto}
\end{equation}
Un cero antes que el polo ($\wL<\wH$): entre ellos la magnitud sube $20\dBdec$ (en total $20\log\alpha$) y la fase se vuelve positiva (\cref{fig:rf-comp}, izquierda). Efectos: más $\PM$ (menos sobrepico) y cruce más a la derecha (más ancho de banda, más rápido). Costo: amplifica el ruido de alta frecuencia en $20\log\alpha$.

\begin{figure}[ht]
\centering
\includegraphics[width=\linewidth]{cap10/f_compensadores_bode}
\caption{Bode del adelanto y del atraso con $\alpha=10$. El adelanto se usa por su fase positiva; el atraso, por su atenuación en alta frecuencia.}
\label{fig:rf-comp}
\end{figure}

\paragraph{¿Dónde está el máximo de fase? Paso a paso}
\begin{enumerate}
  \item La fase del adelanto es $\phi(\omega)=\tan^{-1}\dfrac{\omega}{\wL}-\tan^{-1}\dfrac{\omega}{\wH}$.
  \item Derivando respecto a $\omega$ e igualando a cero, $\dfrac{1/\wL}{1+\omega^2/\wL^2}=\dfrac{1/\wH}{1+\omega^2/\wH^2}$, de donde
  \begin{equation}
    \wM=\sqrt{\wL\,\wH}=\wL\sqrt\alpha .
  \end{equation}
  En escala logarítmica, $\wM$ queda justo en el centro entre el cero y el polo.
  \item Reemplazando, $\tan\phiM=\dfrac{\alpha-1}{2\sqrt\alpha}$, que es lo mismo que
  \begin{equation}
    \sin\phiM=\frac{\alpha-1}{\alpha+1}\qquad\Longleftrightarrow\qquad \alpha=\frac{1+\sin\phiM}{1-\sin\phiM}.
    \label{eq:rf-alpha}
  \end{equation}
  \item En $\wM$ la magnitud está a mitad de camino de su subida: $|\Gcp(\jw_M)|=\sqrt\alpha$, es decir, $+10\log\alpha\dB$.
\end{enumerate}

\paragraph{Rendimientos decrecientes}
La \cref{fig:rf-phimax} muestra $\phiM(\alpha)$. Pasar de $55\grados$ a $65\grados$ exige duplicar $\alpha$ (de $10$ a $20$) y con él el ruido de alta frecuencia. Si se necesitan más de unos $60\grados$, se usan dos etapas de adelanto en cascada o se cambia de estrategia.

\begin{figure}[ht]
\centering
\includegraphics[width=0.78\linewidth]{cap10/f_phimax}
\caption{Fase máxima del adelanto en función de $\alpha$: nunca llega a $90\grados$.}
\label{fig:rf-phimax}
\end{figure}

\paragraph{¿Por qué buscar $|\Gp(\jw_M)|=-10\log\alpha$?}
En $\wM$ el compensador suma $+10\log\alpha\dB$. Si en esa frecuencia la planta vale $-10\log\alpha\dB$, el producto vale $0\dB$: $\wM$ se vuelve el nuevo cruce y el máximo de fase cae exactamente donde se mide el $\PM$. \textbf{¿Y el ángulo de seguridad?} Al correr el cruce de $\wpm$ a $\wM$ (a la derecha) la fase de la planta baja un poco más; $\phiseg$ compensa esa pérdida. Si la pérdida supera a $\phiseg$, se aumenta $\phiseg$ de $5\grados$ en $5\grados$ y se repite.\notamargen{\emph{Caminante, no hay camino, se hace camino al andar.} El diseño en frecuencia es iterativo.}

\begin{receta}[Diseño en adelanto, en siete pasos]
\begin{enumerate}
  \item Escoger $\Kc$ para el error en estado estacionario; anotar $\PMdes$ (o $\zeta$, $\OS$).
  \item Bode de $\Kc\Gp$: medir $\PM$ y $\wpm$.
  \item $\phiM=\PMdes-\PM+\phiseg$, con $\phiseg$ entre $5\grados$ y $10\grados$.
  \item $\alpha=(1+\sin\phiM)/(1-\sin\phiM)$.
  \item Encontrar $\wM$ tal que $20\log|\Kc\Gp(\jw_M)|=-10\log\alpha$.
  \item $\wL=\wM/\sqrt\alpha$ y $\wH=\alpha\,\wL$.
  \item Bode de $\Gc\Gp$: verificar el $\PM$ (y simular). \emph{Verificación estricta:} comprobar que $\phiseg\geq\angle\Gp(\jw_{PM})-\angle\Gp(\jw_M)$; si no, volver al paso 3 con $\phiseg+5\grados$. Si $\phiM\geq90\grados$, el $\PM$ pedido no se alcanza con una sola etapa.
\end{enumerate}
\end{receta}

\begin{ejemplo}[label=ej:rf-adelanto]{Diseño en adelanto}
Diseñe un compensador en adelanto para $\Gp(s)=K/[s(s+2)]$ con error del $5\%$ ante rampa y $\PM=45\grados$.
\solucion
\textbf{Paso 1.} $\ess=1/\kv=0.05\Rightarrow\kv=20$; como $\kv=\lim_{s\to0}s\Gp(s)=K/2$, $K=40$.

\textbf{Paso 2.} Con $K=40$: $\PM=18.0\grados$ en $\wpm=6.17$\,rad/s. Sin compensar, $\OS=60\%$ y $\ts=3.7$\,s.

\textbf{Paso 3.} Con $\phiseg=10\grados$: $\phiM=45-18.0+10=37.0\grados$ (faltan unos $27\grados$; se piden $37\grados$ por la holgura).

\textbf{Paso 4.} $\alpha=(1+\sin37.0\grados)/(1-\sin37.0\grados)=4.03$.

\textbf{Paso 5.} $-10\log\alpha=-6.05\dB$. La magnitud de la planta vale eso en $\wM=8.85$\,rad/s.

\textbf{Paso 6.} $\wL=8.85/\sqrt{4.03}=4.41$ y $\wH=4.03\times4.41=17.76$:
\[
  \Gc(s)=4.03\,\frac{s+4.41}{s+17.76}\qquad(\text{con } K=40 \text{ en la planta}).
\]
\textbf{Paso 7.} $\PM=49.8\grados$ en $8.85$\,rad/s: se cumple con holgura. La verificación estricta también: la planta pierde $5.2\grados$ al correr el cruce de $6.17$ a $8.85$, menos que los $10\grados$ de seguridad. El sobrepico baja de $60\%$ a $22\%$ y $\ts$ de $3.7$ a unos $0.6$\,s (\cref{fig:rf-adelanto-ej}).
\end{ejemplo}

\begin{cuidado}[Una errata del handout anterior]
En la versión anterior de estas notas el texto decía $\wM=14$; era un error tipográfico. Su propio código usaba $8.84$, y los $\wL$ y $\wH$ que reportaba corresponden a $\wM=8.85$.
\end{cuidado}

\begin{figure}[ht]
\centering
\includegraphics[width=\linewidth]{cap10/f_adelanto_ej}
\caption{Ejemplo~\ref{ej:rf-adelanto}: la ``joroba'' de fase del adelanto se ubica en el nuevo cruce y el PM sube de $18.0\grados$ a $49.8\grados$, sin cambiar la ganancia en bajas frecuencias.}
\label{fig:rf-adelanto-ej}
\end{figure}

\codigorepo{cap10_respuesta_frecuencial/Clase4_Compensadores/C4_Ej1_Adelanto.m}{adelanto}{ejemplo \ref{ej:rf-adelanto}}
\codigorepo{cap10_respuesta_frecuencial/Clase4_Compensadores/disenoAdelanto.m}{}{los siete pasos como función}
\video{V7_adelanto_alpha}{$\alpha$ sube de $1$ a $4.03$ con $\wM$ fijo: la joroba crece, el PM sube y el sobrepico baja, con $\kv=20$ todo el tiempo.}

%% ---------------------------------------------------------------------
\subsection{El compensador en atraso: ganancia en bajas, atenuación en el cruce}
\label{sec:rf-atraso}

\begin{equation}
  \Gcp(s)=\frac{1+s/\wH}{1+s/\wL}=\frac{1}{\alpha}\,\frac{s+\wH}{s+\wL},\qquad
  \alpha=\frac{\wH}{\wL}>1 .
  \label{eq:rf-atraso}
\end{equation}
Un polo antes que el cero ($\wL<\wH$). No se usa por su fase (que es negativa) sino por su \emph{atenuación} de $-20\log\alpha$ en alta frecuencia (\cref{fig:rf-comp}, derecha). El truco es poner las esquinas una década por debajo del cruce: allí el atraso ya solo atenúa y ``roba'' apenas unos $5\grados$ de fase.

\begin{receta}[Diseño en atraso]
\begin{enumerate}
  \item Buscar la frecuencia $\omega_{des}$ donde la planta (con el $\Kc$ del error) ya tiene la fase que daría el margen deseado, más una holgura de $5\grados$: $\angle\Gp(\jw_{des})=-180\grados+\PMdes+5\grados$.
  \item Poner el cero una década abajo: $\wH=\omega_{des}/10$.
  \item Escoger $\alpha$ para que $\omega_{des}$ sea el nuevo cruce: $20\log\alpha=20\log|\Gp(\jw_{des})|$; el polo queda en $\wL=\wH/\alpha$.
\end{enumerate}
Es más sencillo que el adelanto (no hay que iterar sobre $\phiM$). El costo: el cruce se mueve a la izquierda y el sistema se vuelve más lento.
\end{receta}

\begin{ejemplo}[label=ej:rf-atraso]{El mismo problema, con atraso}
Repita el diseño del \cref{ej:rf-adelanto} con un compensador en atraso.
\solucion Con $K=40$, la fase vale $-180\grados+45\grados+5\grados=-130\grados$ en $\omega_{des}=1.68$\,rad/s. Allí $|\Gp(j1.68)|=9.13$ ($19.2\dB$), así que $\alpha=9.13$; $\wH=0.168$ y $\wL=0.168/9.13=0.0184$:
\[
  \Gc(s)=\frac{1}{9.13}\,\frac{s+0.168}{s+0.0184}\qquad(\text{con } K=40 \text{ en la planta}),
\]
con $\PM=44.8\grados$. La ganancia en bajas frecuencias no cambió ($\Gcp(0)=1$, $\kv=20$), pero el cruce pasó de $6.17$ a $1.68$\,rad/s (\cref{fig:rf-atraso-ej}).
\end{ejemplo}

\begin{cuidado}[La holgura de $5\grados$ no es opcional]
El handout anterior tomaba $\omega_{des}=2$, es decir, sin la holgura; con ese diseño el $\PM$ queda en unos $40\grados$, por debajo de lo pedido. Los $5\grados$ pagan la fase que el atraso roba en el cruce.
\end{cuidado}

\begin{figure}[ht]
\centering
\includegraphics[width=0.9\linewidth]{cap10/f_atraso_ej}
\caption{Ejemplo~\ref{ej:rf-atraso}: el atraso baja la curva $19.2\dB$ a partir de $\omega\approx0.2$ y conserva la ganancia de bajas frecuencias.}
\label{fig:rf-atraso-ej}
\end{figure}

\codigorepo{cap10_respuesta_frecuencial/Clase4_Compensadores/disenoAtraso.m}{}{los tres pasos como función}

%% ---------------------------------------------------------------------
\subsection{Adelanto contra atraso: el mismo $\kv$, dos personalidades}

\begin{center}\small
\begin{tabular}{@{}lcccccc@{}}
\toprule
 & cruce (rad/s) & $\PM$ & $\OS$ & $\ts$ (s) & $\wbw$ (rad/s) & $u(0^+)$\\
\midrule
sin compensar & 6.17 & $18.0\grados$ & $60\%$ & 3.7 & 9.8 & 1\\
adelanto & 8.85 & $49.8\grados$ & $22\%$ & 0.6 & 14.5 & 4.03\\
atraso & 1.68 & $44.8\grados$ & $27\%$ & 7.1 & 2.8 & 0.11\\
\bottomrule
\end{tabular}
\captionof{table}{Comparación de los dos diseños. $u(0^+)$ es el valor inicial de la señal de control ante un escalón unitario, que es $\Gcp(\infty)$: $\alpha$ para el adelanto y $1/\alpha$ para el atraso.}
\label{tab:rf-comparacion}
\end{center}

\begin{figure}[ht]
\centering
\includegraphics[width=\linewidth]{cap10/f_comparacion}
\caption{Escalón, error ante rampa y esfuerzo de control. Mismo error final ($5\%$): el adelanto llega rápido; el atraso sobrepasa y luego se arrastra hacia $0.05$ (la cola lenta del par polo-cero), pero con mucho menos esfuerzo.}
\label{fig:rf-comparacion}
\end{figure}

\codigorepo{cap10_respuesta_frecuencial/Clase4_Compensadores/C4_Ej3_Comparacion.m}{comparacion}{tabla \ref{tab:rf-comparacion}}
\video{V8_adelanto_vs_atraso}{Siguiendo una rampa: los dos llegan al mismo error, por caminos distintos.}

\begin{anchocompleto}
\centering\small
\begin{tabular}{@{}lll@{}}
\toprule
 & \textbf{Adelanto} & \textbf{Atraso}\\
\midrule
Estructura & $\alpha(s+\wL)/(s+\wH)$, $\alpha>1$ & $(1/\alpha)(s+\wH)/(s+\wL)$, $\alpha>1$\\
Se usa por su\ldots & fase positiva en el cruce & atenuación en el cruce\\
Cruce & se mueve a la derecha & se mueve a la izquierda\\
Velocidad, ancho de banda & aumentan & disminuyen\\
Ruido de alta frecuencia & lo amplifica ($+20\log\alpha$) & lo atenúa\\
Esfuerzo inicial & mayor ($\alpha$) & menor ($1/\alpha$)\\
Límite & unos $60\grados$ por etapa & cola lenta del par polo-cero\\
\bottomrule
\end{tabular}
\captionof{table}{Resumen: adelanto para velocidad; atraso para precisión con poco esfuerzo.}
\label{tab:rf-resumencomp}
\end{anchocompleto}

%% ---------------------------------------------------------------------
\subsection{Atraso-adelanto: cuando hacen falta las dos cosas}

\begin{ejemplo}[label=ej:rf-atrasoadelanto]{Una planta que con el $\kv$ pedido es inestable}
Para $\Gp(s)=K/[s(s+1)(s+2)]$ se pide $\kv=10$, $\PM\geq50\grados$ y $\GM\geq10\dB$.
\solucion $\kv=K/2=10\Rightarrow K=20$, y con esa ganancia el lazo cerrado es inestable: $\PM=-28.1\grados$, $\GM=-10.5\dB$. Un adelanto solo necesitaría unos $85\grados$ de fase; un atraso solo, un cruce muy lento. Se combinan:
\begin{enumerate}
  \item \textbf{Escoger el nuevo cruce}, $\omega_c=1.2$\,rad/s, donde la fase de la planta es $-171.2\grados$.
  \item \textbf{Adelanto centrado en $\omega_c$} con $\alpha=10$ (aporta $\phiM=54.9\grados$): $\wL=1.2/\sqrt{10}=0.38$ y $\wH=1.2\sqrt{10}=3.79$.
  \item \textbf{Atraso que baja la magnitud} en $\omega_c$ lo que sobra para que $\omega_c$ sea el cruce: con el adelanto, $|\Gc\Gp(j1.2)|=23.2\dB$, así que $\beta=10^{23.2/20}=14.5$; el cero va una década abajo ($0.12$) y el polo en $0.12/14.5=0.0083$.
\end{enumerate}
Resultado: $\PM=58.2\grados$ y $\GM=13.2\dB$, con $\kv=10$ intacto (\cref{fig:rf-atrasoadelanto}).
\end{ejemplo}

\begin{figure}[ht]
\centering
\includegraphics[width=\linewidth]{cap10/f_atraso_adelanto}
\caption{Ejemplo~\ref{ej:rf-atrasoadelanto}: el atraso baja la curva en el cruce y el adelanto pone la fase que falta.}
\label{fig:rf-atrasoadelanto}
\end{figure}

\codigorepo{cap10_respuesta_frecuencial/Clase4_Compensadores/C4_Ej4_AtrasoAdelanto.m}{atrasoadelanto}{ejemplo \ref{ej:rf-atrasoadelanto}}

%% ---------------------------------------------------------------------
\subsection{Implementación: de la función de transferencia al circuito}

\begin{figure}[ht]
\centering
\includegraphics[width=0.85\linewidth]{cap10/redes_rc}
\caption{Redes RC pasivas de adelanto y de atraso, con $\alpha=(R_1+R_2)/R_2$. En el adelanto, $\tau=\frac{R_1R_2}{R_1+R_2}C$ y $\alpha\tau=R_1C$; en el atraso, $\tau=R_2C$.}
\label{fig:rf-rc}
\end{figure}

La red de adelanto (\cref{fig:rf-rc}) atenúa en $1/\alpha$: hace falta una etapa de ganancia $\alpha$ (un amplificador operacional) para recuperar el $\kv$. Para el adelanto del \cref{ej:rf-adelanto}, el cero en $\wL=4.41$ da $\alpha\tau=R_1C=0.227$\,s y el polo en $\wH=17.76$ da $\tau=0.0563$\,s; con $C=1\,\mu$F resultan $R_1=227$\,k$\Omega$ y $R_2=R_1/(\alpha-1)=74.9$\,k$\Omega$.

\begin{modelo}[Compensadores.slx]
\cmd{crearModelo\_Compensadores.m} diseña los dos compensadores con \cmd{disenoAdelanto} y \cmd{disenoAtraso} y arma tres lazos en paralelo (sin compensar, adelanto, atraso) con la misma referencia, un selector escalón/rampa y dos osciloscopios: salidas y esfuerzos. Ideal para ver en vivo el error y el sobrepico. La herramienta interactiva \cmd{controlSystemDesigner(Gp)} permite además arrastrar polos y ceros del compensador sobre el Bode y ver el PM en tiempo real.
\end{modelo}

%% ---------------------------------------------------------------------
\section{El PID visto desde la frecuencia}
\label{sec:rf-pid}
%% ---------------------------------------------------------------------

Los compensadores de este capítulo y el PID de los capítulos anteriores son parientes cercanos:
\begin{itemize}
  \item un \textbf{PI}, $\Kc\big(1+\frac{1}{\Ti s}\big)=\Kc\,\frac{\Ti s+1}{\Ti s}$, es un atraso con el polo llevado al origen: ganancia infinita en DC (sube el tipo) a cambio de algo de fase en el cruce;
  \item un \textbf{PD con filtro}, $\Kc\,\frac{1+\Td s}{1+\Td s/N}$, es exactamente un adelanto con $\alpha=N$;
  \item un \textbf{PID} serie con filtro es un atraso-adelanto.
\end{itemize}
Esto permite sintonizar un PID \emph{mirando el Bode}, y entender qué hacen las reglas empíricas.

\paragraph{Ziegler-Nichols leído en frecuencia}
Con la planta del \cref{ej:rf-rangokc} y la sintonización de ZN ($\Kc=73\,326$, $\Ti=0.0943$\,s, $\Td=0.0236$\,s), la malla abierta en la frecuencia crítica vale $\Gc\Gp(\jw_{cr})=-0.6-0.28j$: ZN ``mueve'' el punto que estaba en $-1$ a un lugar fijo, sin importar la planta. El resultado es un lazo agresivo, con $\PM\approx20\grados$ y mucho sobrepico (\cref{fig:rf-pid}).

\paragraph{Moldeo del lazo}
Alternativa: escoger el cruce y el margen y construir el PID para cumplirlos. Con $\omega_c=20$\,rad/s y $\PMdes=50\grados$:
\begin{enumerate}
  \item cero del PI una década abajo del cruce: $\Ti=10/\omega_c=0.5$\,s;
  \item fase que debe aportar el PD filtrado ($N=10$): $-180\grados+50\grados-\big(\angle\Gp+\angle\mathrm{PI}\big)(j20)=37.6\grados$, de donde $\Td=0.0465$\,s;
  \item ganancia para que $|\Gc\Gp(j20)|=1$: $\Kc=33\,427$.
\end{enumerate}
El lazo queda con $\PM=50\grados$ y $\GM=24.4\dB$. La función \cmd{pidtune(Gp,'PIDF',wc,pidtuneOptions('PhaseMargin',50))} hace esto automáticamente.

\begin{figure}[ht]
\centering
\includegraphics[width=0.85\linewidth]{cap10/f_pid_frecuencia}
\caption{Dos PID para la misma planta: Ziegler-Nichols (agresivo) y moldeo del lazo con $\PM=50\grados$.}
\label{fig:rf-pid}
\end{figure}

\codigorepo{cap10_respuesta_frecuencial/Clase4_Compensadores/C4_Ej5_PID_Frecuencia.m}{moldeo}{moldeo del lazo}

\begin{profundiza}[El relé de Åström y Hägglund]
Para medir $\Kcr$ y $\Pcr$ sin llevar la planta al borde de la inestabilidad, Åström y Hägglund propusieron reemplazar el controlador por un relé de amplitud $d$~\citep{AsHa84}. El lazo oscila en la frecuencia donde la fase vale $-180\grados$, con amplitud $a\approx\frac{4d}{\pi\Kcr}$ (aproximación de la función descriptiva), de donde se despeja $\Kcr$. Es el principio de muchos autosintonizadores industriales.
\end{profundiza}

%% ---------------------------------------------------------------------
\section{Para llevar}
\label{sec:rf-cierre}
%% ---------------------------------------------------------------------

\begin{parallevar}
\begin{enumerate}
  \item Un sistema LTI estable responde a $x_0\sin\omega t$ con $x_0|\Gma(\jw)|\sin(\omega t+\angle\Gma(\jw))$. El diagrama de Bode guarda esos dos números para cada $\omega$.
  \item Antes de dibujar, \textbf{forma de Bode}: cada factor como $(s/\omega_i+1)$ y la ganancia $\Kb$ sola. La ganancia escondida es el error más frecuente.
  \item Magnitud: rectas de $\pm20\dBdec$ por factor; error de $3\dB$ en cada esquina real. Fase: $\pm45\grados$/dec entre $\omega_i/10$ y $10\omega_i$. Con polos complejos, \cmd{bode}: la resonancia puede cambiar la conclusión.
  \item La frontera de estabilidad es $\Gma(\jw)=-1$. $\GM$: cuánta ganancia de más se aguanta; $\PM$: cuánto atraso de fase no modelado se aguanta. Regla: $30\grados<\PM<60\grados$, $\GM>6\dB$.
  \item $\PM\leftrightarrow\zeta\leftrightarrow\OS$ (\cref{tab:rf-pmzeta}); el cruce y el ancho de banda dan la velocidad. Las bajas frecuencias dan el tipo y el error.
  \item Diseño: $\Kc$ para el error; luego la forma con $\Gcp(0)=1$. Adelanto para velocidad y margen; atraso para precisión con poco esfuerzo; atraso-adelanto cuando hacen falta las dos cosas. Siempre verificar con \cmd{margin} y en el tiempo.
\end{enumerate}
\end{parallevar}

\begin{profundiza}[Correcciones respecto a la versión anterior del capítulo]
Para quien tenga el handout anterior: (i) adelanto, $\wM=8.85$ (no $14$); (ii) atraso, con la holgura de $5\grados$, $\omega_{des}=1.68$ y $\PM=44.8\grados$; (iii) rango de $\Kc$, $\GM=+40\dB$ (no $-40\dB$); (iv) $\wgm$ asintótico del ejemplo $10\,000/[\cdots]$, $31.6$ (no $50$); (v) ajuste de ganancia, $\Kc$ exacto $33\,056$; (vi) el pico de resonancia se llama ahora $\Mr$ (antes $M_P$) para no confundirlo con el sobrepico. Todos los números del capítulo se verificaron de forma independiente (\texttt{herramientas/cap10\_numeros.py} en el repositorio y \cmd{verificarCap10.m} en \matlab).
\end{profundiza}

%% ---------------------------------------------------------------------
\section{Problemas}
\label{sec:rf-problemas}
%% ---------------------------------------------------------------------

Los problemas marcados con \etiquetaM{} requieren \matlab/\simulink; los marcados con \etiquetaEstrella{} son de profundización. Al final hay soluciones de los primeros.

\subsection*{Diagramas asintóticos y respuesta senoidal}

\begin{problema}[asintótico]\label{prob:rf-1}
Dibuje el diagrama asintótico de magnitud y fase de $\Gma(s)=\dfrac{100(s+1)}{s(s+10)}$. Marque el valor en $\omega=1$ y compare con \cmd{bode} en las esquinas.
\end{problema}

\begin{problema}[asintótico]\label{prob:rf-2}
Ídem para $\Gma(s)=\dfrac{50}{(s+2)(s+5)}$.
\end{problema}

\begin{problema}[respuesta senoidal]\label{prob:rf-3}
Para $\Gma(s)=\dfrac{10}{s+10}$ y $x(t)=2\sin(10t)$, calcule $z_{ss}(t)$ a mano y verifique con \cmd{lsim}.
\end{problema}

\begin{problema}[Simulink]\label{prob:rf-4}
\etiquetaM{} Ejecute \cmd{crearModelo\_RespuestaSenoidal.m} con $G(s)=100/(s^2+4s+100)$. Mida amplitud y desfase para $\omega=1$, $10$ y $50$ y ubique los tres puntos sobre el Bode de \matlab. ¿En cuál de los tres la salida es mayor que la entrada? ¿Por qué?
\end{problema}

\begin{problema}[repaso]\label{prob:rf-5}
Sea $\Gma(s)=\dfrac{20(s+2)}{s(s+20)}$. (a) Escríbala en forma de Bode. (b) Dibuje el diagrama asintótico de magnitud y fase. (c) Compare con los valores reales en $\omega=2$ y $\omega=20$. (d) ¿Cuál es el tipo del sistema y cuánto vale $\kv$?
\end{problema}

\begin{problema}[repaso]\label{prob:rf-6}
Sea $G(s)=\dfrac{50}{(s+5)(s+10)}$ con $x(t)=3\sin(5t)$. (a) Calcule $z_{ss}(t)$ exacta. (b) Estímela solo con el diagrama asintótico y compare. (c) ¿A qué frecuencia la salida es diez veces más pequeña que la entrada (asintótico y exacto)?
\end{problema}

\begin{problema}[el error de la asíntota]\label{prob:rf-7}
Para las mallas abiertas (i) $\dfrac{200}{(s+1)(s+20)}$, (ii) $\dfrac{50(s+2)}{s(s+10)(s+50)}$ y (iii) $\dfrac{5(s+1)}{s^2(s+5)}$:
(a) escríbalas en forma de Bode e identifique $\Kb$, el número de integradores y las esquinas;
(b) justifique a partir de la expresión de la magnitud las reglas asintóticas de un polo real simple;
(c) trace los tres diagramas indicando la pendiente de cada tramo;
(d) calcule el valor exacto de magnitud y fase en cada esquina y verifique que el error de magnitud en una esquina simple es siempre el mismo;
(e) para (ii), determine analíticamente dónde es máximo el error entre la asíntota y la curva real en el tramo $2<\omega<10$ y cuantifíquelo;
(f) superponga los diagramas sobre los de \cmd{bode} y comente dónde la aproximación es aceptable.
\end{problema}

\subsection*{Estabilidad y márgenes}

\begin{problema}[identificación]\label{prob:rf-8}
Un sistema de fase mínima y ganancia positiva tiene un diagrama asintótico de magnitud con pendiente de $-20\dBdec$ en bajas frecuencias, pasando por $40\dB$ en $\omega=0.1$; quiebre a $0\dBdec$ en $\omega=2$; quiebre a $-40\dBdec$ en $\omega=20$, y quiebre a $-20\dBdec$ en $\omega=200$. (a) Reconstruya $\Gma(s)$, justificando cada factor. (b) Trace la fase y determine $\PM$ y $\GM$. (c) Al medir la fase se encuentra que en alta frecuencia tiende a $-270\grados$ en lugar del valor predicho. ¿Qué modificación de $\Gma(s)$ es compatible con la misma magnitud y explica esa fase? (d) Demuestre que ambas funciones tienen diagramas de magnitud idénticos. (e) Compare sus fases en $\omega=20$, $200$ y $2000$\,rad/s y concluya cuál es estable en lazo cerrado con retroalimentación unitaria. (f) Enuncie qué información \emph{no} contiene el diagrama de magnitud y bajo qué condición la magnitud determina la fase.
\end{problema}

\begin{problema}[márgenes y rango de ganancia]\label{prob:rf-9}
Para $\Gma(s)=\dfrac{10\,000}{(s+1)(s+10)(s+100)}$:
(a) trace el asintótico de magnitud y halle $\wpm$;
(b) trace el asintótico de fase con sus puntos de quiebre y halle $\wgm$ y el $\GM$;
(c) compare los cuatro valores con \cmd{margin} y calcule el error porcentual de cada uno;
(d) con $\Gp=\dfrac{1}{(s+1)(s+10)(s+100)}$ y $\Gc=\Kc$, halle el rango de $\Kc$ con Routh-Hurwitz;
(e) obtenga $\Kcr$ del $\GM$ leído en el Bode, compare y explique por qué el error de la lectura asintótica va en la dirección conveniente para el diseño;
(f) repita (a)--(e) para $\Gp=\dfrac{1}{s(s+2)(s+10)}$, obteniendo además la frecuencia crítica por sustitución $s=\jw$; con $\Kcr$ y $\Pcr$ sintonice un PID por Ziegler-Nichols y determine el margen de fase resultante. ¿Es aceptable?
\end{problema}

\begin{problema}[presupuesto de fase]\label{prob:rf-10}
Al lazo del problema~\ref{prob:rf-9}(a) se le agrega el sensor $\Hs(s)=20/(s+20)$. (a) Estime el nuevo $\PM$ con $\PM-\tan^{-1}(\wpm/20)$. (b) Calcúlelo con \cmd{margin}. (c) ¿Por qué la estimación es pesimista? (d) \etiquetaM{} Simule el escalón con y sin sensor y compare los sobrepicos.
\end{problema}

\subsection*{Lazo cerrado, tipo y ajuste de ganancia}

\begin{problema}[PM y transitorio]\label{prob:rf-11}
Para $\Gma(s)=\dfrac{K}{s(s+5)}$ con retroalimentación unitaria:
(a) con $K=100$, identifique $\wn$ y $\zeta$ del lazo cerrado por comparación directa;
(b) obtenga lo mismo por la vía frecuencial ($\omega_1$, $\PM$, ecuación~\eqref{eq:rf-PMzeta});
(c) calcule $\OS$, $\ts$ y $\tp$ y valide simulando;
(d) calcule $\wbw$ con \eqref{eq:rf-BW} y verifíquelo en el Bode del lazo cerrado;
(e) ¿hay pico de resonancia? Si lo hay, calcule $\Mr$ y $\wres$;
(f) halle el $K$ que da $\PM=60\grados$ y compare con $K=100$ en términos de $\wbw$: ¿cuál es más rápido y cuál más amortiguado?
\end{problema}

\begin{problema}[tipo y error]\label{prob:rf-12}
Para (i) $\dfrac{20}{(s+2)(s+5)}$, (ii) $\dfrac{30}{s(s+3)(s+10)}$ y (iii) $\dfrac{400(s+2)}{s^2(s+8)}$ con retroalimentación unitaria:
(a) determine el tipo mirando solo la pendiente de baja frecuencia;
(b) lea gráficamente la constante de error;
(c) verifíquela analíticamente;
(d) calcule el error ante escalón, rampa y parábola, verificando primero la estabilidad;
(e) uno de los tres no admite el cálculo de (d): ¿cuál, por qué y qué rasgo de su Bode lo anticipa?
\end{problema}

\begin{problema}[ajuste de ganancia]\label{prob:rf-13}
Con $\Gc=\Kc$ y $\Gp=\dfrac{1}{s(s+20)(s+100)}$:
(a) halle $\zeta$ y el $\PM$ para $\OS=10\%$;
(b) en el Bode de $2000\,\Gp$ localice la frecuencia con la fase objetivo y lea la magnitud;
(c) determine $\Kc$ y $\kv$;
(d) \etiquetaM{} adapte \cmd{C3\_Ej3\_AjusteGanancia.m}, simule y mida el $\OS$ real;
(e) explique por qué no coincide exactamente con el especificado;
(f) calcule el $\GM$ del sistema ajustado y compárelo con el criterio $\GM>6\dB$.
\end{problema}

\subsection*{Compensadores}

\begin{problema}[adelanto]\label{prob:rf-14}
Diseñe un compensador en adelanto para $\Gp(s)=\dfrac{100K}{s(s+36)(s+100)}$ con $\kv=40$ y $\OS=20\%$. Use $\phiseg=10\grados$ y aplique la verificación estricta; si el $\PM$ no se cumple, itere. \etiquetaM{} Verifique con \cmd{disenoAdelanto(Gp,48.1,10,true)}.
\end{problema}

\begin{problema}[atraso y comparación]\label{prob:rf-15}
Para la planta del problema anterior diseñe un compensador en atraso con las mismas especificaciones (\cmd{disenoAtraso(Gp,48.1)}) y construya una tabla como la \cref{tab:rf-comparacion}. ¿En qué aplicación preferiría cada uno?
\end{problema}

\begin{problema}[implementación]\label{prob:rf-16}
Calcule $R_1$ y $R_2$ de la red de adelanto del problema~\ref{prob:rf-14} con $C=1\,\mu$F. ¿Qué ganancia debe tener el amplificador que la acompaña?
\end{problema}

\begin{problema}[límites del adelanto]\label{prob:rf-17}
\etiquetaEstrella{} Con $\Gp=K/[s(s+2)]$ exija $\kv\geq200$ y $\PM\geq45\grados$. (a) Halle $K$ y el $\PM$ sin compensar. (b) ¿Qué $\alpha$ necesitaría un adelanto de una sola etapa? (c) Repita con $\PM\geq65\grados$ y evalúe si es realizable (un adelanto de primer orden no aporta $90\grados$ y un $\alpha$ muy grande amplifica el ruido). (d) Si no lo es, diseñe una alternativa (dos etapas o atraso-adelanto) y verifique que cumple.
\end{problema}

%% ---------------------------------------------------------------------
\section{Soluciones seleccionadas}
%% ---------------------------------------------------------------------

\paragraph{Problema~\ref{prob:rf-1}}
$\Gma(s)=\dfrac{10(s+1)}{s(s/10+1)}$: $\Kb=10$, un integrador, un cero en $1$ y un polo en $10$. La magnitud baja a $-20\dBdec$ y pasa por $20\dB$ en $\omega=1$; entre $1$ y $10$ queda plana en $20\dB$; después cae a $-20\dBdec$. Fase asintótica: $-90\grados$ hasta $0.1$, sube hasta $-45\grados$ en $\omega=1$, plana en $-45\grados$ hasta $10$, y vuelve a $-90\grados$ en $100$. Valores reales: $22.97\dB$ en $\omega=1$ y $17.03\dB$ en $\omega=10$ (la asíntota dice $20\dB$ en ambos: error de $\pm3\dB$), y $-50.7\grados$ en ambas esquinas.

\paragraph{Problema~\ref{prob:rf-2}}
$\Gma(s)=\dfrac{5}{(s/2+1)(s/5+1)}$: $13.98\dB$ hasta $\omega=2$, $-20\dBdec$ entre $2$ y $5$, $-40\dBdec$ después; la fase va de $0\grados$ a $-180\grados$. Valor real en $\omega=2$: $10.33\dB$ y $-66.8\grados$.

\paragraph{Problema~\ref{prob:rf-3}}
$|\Gma(j10)|=10/|10+j10|=0.707$ y $\phi=-45\grados$: $z_{ss}(t)=1.414\sin(10t-45\grados)$.

\paragraph{Problema~\ref{prob:rf-4}}
$\omega=1$: $1.01$ ($+0.08\dB$), $-2.3\grados$; $\omega=10$: $2.5$ ($+7.96\dB$), $-90\grados$; $\omega=50$: $0.042$ ($-27.6\dB$), $-175.2\grados$. La salida es mayor que la entrada en $\omega=10$, la frecuencia natural ($\zeta=0.2$).

\paragraph{Problema~\ref{prob:rf-5}}
(a) $20(s+2)=40(s/2+1)$ y $s(s+20)=20\,s\,(s/20+1)$, luego $\Gma(s)=\dfrac{2\,(s/2+1)}{s\,(s/20+1)}$, $\Kb=2$ ($6.02\dB$).
(b) Magnitud asintótica: $6.02-20\log\omega$ para $\omega<2$; $0\dB$ entre $2$ y $20$ (el integrador y el cero se cancelan); $-20\log(\omega/20)$ después. Fase: $-90\grados$ hasta $0.2$; entre $2$ y $20$ el cero y el polo se compensan y la fase queda plana en $-45\grados$; vuelve a $-90\grados$ en $200$.
(c) $|\Gma(j2)|=+2.97\dB$ y $|\Gma(j20)|=-2.97\dB$ (la asíntota dice $0$: el error de cada esquina); fase real $-50.7\grados$ en ambas (asintótica: $-45\grados$).
(d) Tipo 1, $\kv=\lim_{s\to0}s\Gma(s)=2$; error ante rampa unitaria $1/\kv=0.5$.

\paragraph{Problema~\ref{prob:rf-6}}
(a) $|G(j5)|=\dfrac{50}{7.071\cdot11.18}=0.632$ ($-3.98\dB$), $\phi=-45\grados-26.6\grados=-71.6\grados$, luego $z_{ss}(t)=1.90\sin(5t-71.6\grados)$.
(b) En forma de Bode $G=1/[(s/5+1)(s/10+1)]$: en $\omega=5$ la asíntota vale $0\dB$ y la fase asintótica $-45\grados-45\grados\log_{10}5=-76.5\grados$, así que $z_{ss}\approx3\sin(5t-76.5\grados)$. La asíntota se equivoca en $3.98\dB$ (un factor $1.58$): sirve para estimar; para calcular se usa $G(\jw)$.
(c) Asintótico: $-6.02-40\log(\omega/10)=-20\Rightarrow\omega\approx22.4$\,rad/s. Exacto: $\omega^4+125\omega^2-247\,500=0\Rightarrow\omega=20.95$\,rad/s.

\codigorepo{cap10_respuesta_frecuencial/Clase1_RespuestaFrecuencial/C1_Ej4_Repaso.m}{}{verificación de los problemas \ref{prob:rf-5} y \ref{prob:rf-6}}

\respuestaspausas

